% Vietnamese translation for OI Public Library
% Source-PDF: 英文资料 ENGLISH MATERIALS/SatoNT.pdf
\documentclass[11pt]{article}
\usepackage{oipl-vn}

\DeclareMathOperator{\lcm}{lcm}
\DeclareMathOperator{\ord}{ord}
\DeclareMathOperator{\ind}{ind}
\newcommand{\leg}[2]{\left(\frac{#1}{#2}\right)}
\newcommand{\Z}{\mathbb{Z}}
\newcommand{\ds}{\displaystyle}

\title{Lý thuyết số}
\author{Naoki Sato}
\date{}

\begin{document}
\maketitle

\section*{Lời nói đầu}

Bộ ghi chú về lý thuyết số này ban đầu được viết năm 1995 cho học sinh ở
trình độ IMO. Tài liệu trình bày những kiến thức nền tảng mà một học sinh IMO
nên nắm vững. Văn bản này được dùng như tài liệu tham khảo, không thay thế mà
bổ sung cho sách giáo khoa lý thuyết số; một vài tài liệu như vậy được liệt
kê ở cuối. Chứng minh được đưa ra khi thích hợp, hoặc khi chúng minh họa một
ý tưởng quan trọng. Các bài toán được chọn lọc từ nhiều nguồn, trong đó có
nhiều bài từ các kỳ thi và olympiad, và nhìn chung khá khó. Tác giả hoan nghênh
mọi chỉnh sửa và góp ý.

\section{Tính chia hết}

Với các số nguyên \(a\) và \(b\), ta nói \(a\) chia hết \(b\), hoặc \(a\) là
ước của \(b\), hoặc \(b\) là bội của \(a\), nếu tồn tại số nguyên \(c\) sao
cho \(b=ca\). Ta ký hiệu điều này là \(a\mid b\). Nếu không, ta viết
\(a\nmid b\). Một số nguyên dương \(p\) là số nguyên tố nếu các ước dương duy
nhất của \(p\) là \(1\) và \(p\). Nếu \(p^k\mid a\) nhưng \(p^{k+1}\nmid a\),
trong đó \(p\) là số nguyên tố, tức \(p^k\) là lũy thừa cao nhất của \(p\)
chia hết \(a\), ta viết \(p^k\|a\).

\paragraph{Các sự kiện hữu ích.}
\begin{itemize}
\item Nếu \(a,b>0\) và \(a\mid b\), thì \(a\le b\).
\item Nếu \(a\mid b_1,a\mid b_2,\ldots,a\mid b_n\), thì với mọi số nguyên
\(c_1,c_2,\ldots,c_n\),
\[
  a\mid \sum_{i=1}^n b_i c_i.
\]
\end{itemize}

\paragraph{Định lý 1.1 (thuật toán chia).}
Với mọi số nguyên dương \(a\) và số nguyên \(b\), tồn tại duy nhất các số
nguyên \(q,r\) sao cho \(b=qa+r\) và \(0\le r<a\), trong đó \(r=0\) khi và chỉ
khi \(a\mid b\).

\paragraph{Định lý 1.2 (định lý cơ bản của số học).}
Mọi số nguyên lớn hơn \(1\) đều viết được một cách duy nhất dưới dạng
\[
  p_1^{e_1}p_2^{e_2}\cdots p_k^{e_k},
\]
trong đó các \(p_i\) là các số nguyên tố phân biệt và các \(e_i\) là số nguyên
dương.

\paragraph{Định lý 1.3 (Euclid).}
Có vô hạn số nguyên tố.

\paragraph{Chứng minh.}
Giả sử chỉ có hữu hạn số nguyên tố, là \(p_1,p_2,\ldots,p_n\). Đặt
\(N=p_1p_2\cdots p_n+1\). Theo định lý cơ bản của số học, \(N\) chia hết cho
một số nguyên tố \(p\). Theo giả thiết, \(p\) phải nằm trong các \(p_i\), nhưng
rõ ràng \(N\) không chia hết cho bất kỳ \(p_i\) nào, mâu thuẫn.

\paragraph{Ví dụ 1.1.}
Cho \(x,y\) là các số nguyên. Chứng minh rằng \(2x+3y\) chia hết cho \(17\)
khi và chỉ khi \(9x+5y\) chia hết cho \(17\).

\paragraph{Lời giải.}
Nếu \(17\mid(2x+3y)\), thì \(17\mid 13(2x+3y)\), tức
\(17\mid(26x+39y)\), suy ra \(17\mid(9x+5y)\). Ngược lại, nếu
\(17\mid(9x+5y)\), thì \(17\mid4(9x+5y)\), tức \(17\mid(36x+20y)\), suy ra
\(17\mid(2x+3y)\).

\paragraph{Ví dụ 1.2.}
Tìm tất cả số nguyên dương \(d\) sao cho với một số nguyên \(n\) nào đó, \(d\)
chia hết cả \(n^2+1\) và \((n+1)^2+1\).

\paragraph{Lời giải.}
Giả sử \(d\mid(n^2+1)\) và \(d\mid((n+1)^2+1)=n^2+2n+2\). Khi đó
\[
 d\mid(n^2+2n+2)-(n^2+1)=2n+1.
\]
Suy ra \(d\mid4n^2+4n+1\), nên
\[
 d\mid4(n^2+2n+2)-(4n^2+4n+1)=4n+7.
\]
Do đó \(d\mid(4n+7)-2(2n+1)=5\). Vậy \(d\) chỉ có thể là \(1\) hoặc \(5\).
Lấy \(n=2\) cho thấy cả hai giá trị đều đạt được.

\paragraph{Ví dụ 1.3.}
Giả sử \(a_1,a_2,\ldots,a_{2n}\) là các số nguyên phân biệt sao cho phương
trình
\[
  (x-a_1)(x-a_2)\cdots(x-a_{2n})-(-1)^n(n!)^2=0
\]
có nghiệm nguyên \(r\). Chứng minh rằng
\[
  r=\frac{a_1+a_2+\cdots+a_{2n}}{2n}.
\]
(Danh sách chọn lọc IMO 1984)

\paragraph{Lời giải.}
Rõ ràng \(r\ne a_i\) với mọi \(i\), và các số \(r-a_i\) là \(2n\) số nguyên
phân biệt. Vì vậy
\[
 |(r-a_1)\cdots(r-a_{2n})|
 \ge |1\cdot2\cdots n\cdot(-1)\cdot(-2)\cdots(-n)|=(n!)^2,
\]
dấu bằng xảy ra khi và chỉ khi
\[
 \{r-a_1,\ldots,r-a_{2n}\}=\{1,2,\ldots,n,-1,-2,\ldots,-n\}.
\]
Do phương trình đã cho buộc dấu bằng xảy ra, ta có
\[
\begin{aligned}
 0&=(r-a_1)+\cdots+(r-a_{2n})\\
  &=2nr-(a_1+\cdots+a_{2n}),
\end{aligned}
\]
từ đó suy ra kết luận.

\paragraph{Ví dụ 1.4.}
Cho \(0<a_1<a_2<\cdots<a_{mn+1}\) là \(mn+1\) số nguyên. Chứng minh rằng ta có
thể chọn hoặc \(m+1\) số trong đó sao cho không số nào chia hết số nào, hoặc
\(n+1\) số trong đó sao cho mỗi số chia hết số đứng sau.
(Putnam 1966)

\paragraph{Lời giải.}
Với mỗi \(i\), \(1\le i\le mn+1\), gọi \(n_i\) là độ dài của dãy dài nhất bắt
đầu từ \(a_i\), trong các số \(a_i,a_{i+1},\ldots,a_{mn+1}\), sao cho mỗi số
chia hết số tiếp theo. Nếu có \(n_i>n\), bài toán xong. Ngược lại, theo nguyên
lý Dirichlet, có ít nhất \(m+1\) giá trị \(n_i\) bằng nhau. Các số \(a_i\)
tương ứng với các giá trị ấy không thể chia hết lẫn nhau.

\paragraph{Các sự kiện hữu ích.}
\begin{itemize}
\item Tiên đề Bertrand: với mọi số nguyên dương \(n\), tồn tại số nguyên tố
\(p\) sao cho \(n\le p\le2n\).
\item Bổ đề Gauss: nếu một đa thức hệ số nguyên phân tích được thành hai đa
thức hệ số hữu tỉ, thì nó cũng phân tích được thành hai đa thức hệ số nguyên.
\end{itemize}

\paragraph{Bài tập.}
\begin{enumerate}
\item Cho \(a,b\) là các số nguyên dương sao cho
\[
 a\mid b^2,\quad b^2\mid a^3,\quad a^3\mid b^4,\quad b^4\mid a^5,\ldots.
\]
Chứng minh \(a=b\).
\item Cho \(a,b,c\) là ba số nguyên phân biệt, và \(P\) là đa thức có mọi hệ
số nguyên. Chứng minh rằng không thể có \(P(a)=b\), \(P(b)=c\), và \(P(c)=a\).
(USAMO 1974)
\item Chứng minh rằng nếu \(a,b\) là các số nguyên dương, thì
\[
 \left(a+\frac12\right)^n+\left(b+\frac12\right)^n
\]
là số nguyên chỉ với hữu hạn số nguyên dương \(n\). (A Problem Seminar, D.J.
Newman)
\item Với số nguyên dương \(n\), gọi \(r(n)\) là tổng các số dư khi chia \(n\)
cho \(1,2,\ldots,n\). Chứng minh rằng \(r(k)=r(k-1)\) với vô hạn số nguyên
dương \(k\). (Kürschák 1981)
\item Chứng minh rằng với mọi số nguyên dương \(n\),
\[
 0<\sum_{k=1}^n\frac{g(k)}{k}-\frac{2n}{3}<\frac23,
\]
trong đó \(g(k)\) là ước lẻ lớn nhất của \(k\). (Olympiad Toán Áo 1973)
\item Cho \(d\) là số nguyên dương, và \(S\) là tập tất cả số nguyên dương có
dạng \(x^2+dy^2\), trong đó \(x,y\) là các số nguyên không âm.
\begin{enumerate}[label=(\alph*)]
\item Chứng minh rằng nếu \(a\in S\) và \(b\in S\), thì \(ab\in S\).
\item Chứng minh rằng nếu \(a\in S\), \(p\in S\), \(p\) là số nguyên tố và
\(p\mid a\), thì \(a/p\in S\).
\item Giả sử phương trình \(x^2+dy^2=p\) có nghiệm nguyên không âm \(x,y\),
trong đó \(p\) là số nguyên tố cho trước. Chứng minh rằng nếu \(d\ge2\) thì
nghiệm là duy nhất, còn nếu \(d=1\) thì có đúng hai nghiệm.
\end{enumerate}
\end{enumerate}

\section{Ước chung lớn nhất và bội chung nhỏ nhất}

Ước chung lớn nhất của hai số nguyên dương \(a,b\) là số nguyên dương lớn nhất
chia hết cả \(a\) và \(b\), ký hiệu \(\gcd(a,b)\). Tương tự, bội chung nhỏ
nhất của \(a,b\) là số nguyên dương nhỏ nhất là bội của cả \(a\) và \(b\), ký
hiệu \(\lcm(a,b)\). Ta nói \(a,b\) nguyên tố cùng nhau nếu \(\gcd(a,b)=1\).
Với các số nguyên \(a_1,\ldots,a_n\), \(\gcd(a_1,\ldots,a_n)\) là số nguyên
dương lớn nhất chia hết tất cả chúng, và \(\lcm(a_1,\ldots,a_n)\) được định
nghĩa tương tự.

\paragraph{Các sự kiện hữu ích.}
\begin{itemize}
\item Với mọi \(a,b\), \(\gcd(a,b)\lcm(a,b)=ab\).
\item Với mọi \(a,b,m\), \(\gcd(ma,mb)=m\gcd(a,b)\) và
\(\lcm(ma,mb)=m\lcm(a,b)\).
\item Nếu \(d\mid\gcd(a,b)\), thì
\[
 \gcd\left(\frac ad,\frac bd\right)=\frac{\gcd(a,b)}d.
\]
Đặc biệt, nếu \(d=\gcd(a,b)\), thì \(\gcd(a/d,b/d)=1\).
\item Nếu \(a\mid bc\) và \(\gcd(a,c)=1\), thì \(a\mid b\).
\item Nếu \(d\) là ước chung dương của \(a,b\) và mọi ước chung \(d'\) của
\(a,b\) đều chia hết \(d\), thì \(d=\gcd(a,b)\).
\item Nếu \(a_1a_2\cdots a_n\) là một lũy thừa bậc \(k\) hoàn hảo và các
\(a_i\) đôi một nguyên tố cùng nhau, thì từng \(a_i\) cũng là lũy thừa bậc
\(k\) hoàn hảo.
\item Hai số nguyên liên tiếp bất kỳ nguyên tố cùng nhau.
\end{itemize}

\paragraph{Ví dụ 2.1.}
Chứng minh rằng với mọi số nguyên dương \(N\), tồn tại một bội của \(N\) chỉ
gồm các chữ số \(1\) và \(0\). Hơn nữa, nếu \(N\) nguyên tố cùng nhau với
\(10\), thì tồn tại một bội chỉ gồm các chữ số \(1\).

\paragraph{Lời giải.}
Xét \(N+1\) số \(1,11,111,\ldots,\underbrace{11\cdots1}_{N+1\text{ chữ số}}\).
Khi chia cho \(N\), chúng cho \(N+1\) số dư. Theo nguyên lý Dirichlet, hai số
dư bằng nhau, nên hiệu của hai số tương ứng, có dạng \(11\cdots100\cdots0\),
chia hết cho \(N\). Nếu \(N\) nguyên tố cùng nhau với \(10\), ta chia bỏ mọi
lũy thừa của \(10\) để thu được một số dạng \(11\cdots1\) vẫn chia hết cho
\(N\).

\paragraph{Định lý 2.1.}
Với mọi số nguyên dương \(a,b\), tồn tại các số nguyên \(x,y\) sao cho
\[
 ax+by=\gcd(a,b).
\]
Hơn nữa, khi \(x,y\) chạy qua mọi số nguyên, \(ax+by\) nhận đúng các bội của
\(\gcd(a,b)\).

\paragraph{Chứng minh.}
Gọi \(S\) là tập các số nguyên có dạng \(ax+by\), và gọi \(d\) là phần tử
dương nhỏ nhất của \(S\). Theo thuật toán chia, tồn tại \(q,r\) sao cho
\(a=qd+r\), \(0\le r<d\). Khi đó
\[
 r=a-qd=a-q(ax+by)=(1-qx)a-(qy)b,
\]
nên \(r\in S\). Vì \(r<d\), suy ra \(r=0\), do đó \(d\mid a\). Tương tự,
\(d\mid b\), nên \(d\mid\gcd(a,b)\). Mặt khác, \(\gcd(a,b)\) chia mọi phần tử
của \(S\), đặc biệt chia \(d\). Vậy \(d=\gcd(a,b)\). Phần còn lại suy ra ngay.

\paragraph{Hệ quả 2.2.}
Hai số nguyên dương \(a,b\) nguyên tố cùng nhau khi và chỉ khi tồn tại các số
nguyên \(x,y\) sao cho \(ax+by=1\).

\paragraph{Hệ quả 2.3.}
Với mọi số nguyên dương \(a_1,\ldots,a_n\), tồn tại các số nguyên
\(x_1,\ldots,x_n\) sao cho
\[
 a_1x_1+a_2x_2+\cdots+a_nx_n=\gcd(a_1,a_2,\ldots,a_n).
\]

\paragraph{Hệ quả 2.4.}
Cho \(a,b\) là các số nguyên dương và \(n\) là số nguyên. Phương trình
\[
 ax+by=n
\]
có nghiệm nguyên \(x,y\) khi và chỉ khi \(\gcd(a,b)\mid n\). Nếu có nghiệm,
mọi nghiệm có dạng
\[
 (x,y)=\left(x_0+t\frac bd,\;y_0-t\frac ad\right),
\]
trong đó \(d=\gcd(a,b)\), \((x_0,y_0)\) là một nghiệm cụ thể, và \(t\) là số
nguyên.

\paragraph{Chứng minh.}
Phần đầu theo Định lý 2.1. Với phần sau, đặt \(d=\gcd(a,b)\), và giả sử
\((x_0,y_0)\) là một nghiệm. Nếu \(ax+by=n\), thì
\[
 a(x-x_0)+b(y-y_0)=0,
\]
nên \(a(x-x_0)=b(y_0-y)\), hay
\[
 (x-x_0)\frac ad=(y_0-y)\frac bd.
\]
Vì \(a/d\) và \(b/d\) nguyên tố cùng nhau, \(b/d\mid x-x_0\) và
\(a/d\mid y_0-y\). Đặt \(x-x_0=tb/d\), \(y_0-y=ta/d\), ta được dạng nghiệm đã
nêu.

\paragraph{Ví dụ 2.2.}
Chứng minh rằng phân số
\[
 \frac{21n+4}{14n+3}
\]
tối giản với mọi số nguyên dương \(n\). (IMO 1959)

\paragraph{Lời giải.}
Với mọi \(n\),
\[
 3(14n+3)-2(21n+4)=1,
\]
nên tử và mẫu nguyên tố cùng nhau.

\paragraph{Ví dụ 2.3.}
Với mọi số nguyên dương \(n\), đặt \(T_n=2^{2^n}+1\). Chứng minh rằng nếu
\(m\ne n\), thì \(T_m\) và \(T_n\) nguyên tố cùng nhau.

\paragraph{Lời giải.}
Ta có
\[
\begin{aligned}
 T_n-2&=2^{2^n}-1=(T_{n-1}-1)^2-1\\
 &=T_{n-1}(T_{n-1}-2)=T_{n-1}T_{n-2}(T_{n-2}-2)\\
 &=\cdots=T_{n-1}T_{n-2}\cdots T_1T_0.
\end{aligned}
\]
Do đó mọi ước chung của \(T_m\) và \(T_n\) phải chia \(2\). Nhưng mọi \(T_n\)
đều lẻ, nên chúng nguyên tố cùng nhau. Kết quả này lập tức cho một chứng minh
khác rằng có vô hạn số nguyên tố.

\paragraph{Thuật toán Euclid.}
Bằng cách dùng lặp lại thuật toán chia, ta tìm được \(\gcd(a,b)\) mà không cần
phân tích \(a,b\), đồng thời tìm được \(x,y\) trong Định lý 2.1. Ví dụ, với
\(a=329\), \(b=182\),
\[
 329=1\cdot182+147,\quad
 182=1\cdot147+35,\quad
 147=4\cdot35+7,\quad
 35=5\cdot7.
\]
Ước cuối cùng khác không là \(\gcd(329,182)=7\). Truy ngược,
\[
\begin{aligned}
7&=147-4\cdot35=147-4(182-147)\\
 &=5\cdot147-4\cdot182=5(329-182)-4\cdot182\\
 &=5\cdot329-9\cdot182.
\end{aligned}
\]
Thuật toán Euclid cũng áp dụng cho đa thức.

\paragraph{Ví dụ 2.4.}
Cho \(n\) là số nguyên dương và \(S\) là tập con gồm \(n+1\) phần tử của
\(\{1,2,\ldots,2n\}\). Chứng minh rằng:
\begin{enumerate}[label=(\alph*)]
\item tồn tại hai phần tử của \(S\) nguyên tố cùng nhau;
\item tồn tại hai phần tử của \(S\) sao cho một số chia hết số kia.
\end{enumerate}

\paragraph{Lời giải.}
(a) Trong \(S\) phải có hai số liên tiếp, nên chúng nguyên tố cùng nhau.

(b) Xét ước lẻ lớn nhất của mỗi trong \(n+1\) phần tử của \(S\). Mỗi ước nằm
trong \(n\) số lẻ \(1,3,\ldots,2n-1\). Theo nguyên lý Dirichlet, có hai số có
cùng ước lẻ lớn nhất; chúng chỉ khác nhau bởi một lũy thừa của \(2\), nên một
số chia hết số kia.

\paragraph{Ví dụ 2.5.}
Các số nguyên dương \(a_1,\ldots,a_n\) đều nhỏ hơn \(1000\), và
\(\lcm(a_i,a_j)>1000\) với mọi \(i\ne j\). Chứng minh rằng
\[
 \sum_{i=1}^n\frac1{a_i}<2.
\]
(Olympiad Toán Nga 1951)

\paragraph{Lời giải.}
Nếu \(\frac{1000}{m+1}<a\le\frac{1000}{m}\), thì các bội
\(a,2a,\ldots,ma\) không vượt quá \(1000\). Gọi \(k_1\) là số \(a_i\) trong
\((1000/2,1000]\), \(k_2\) trong \((1000/3,1000/2]\), v.v. Khi đó có
\(k_1+2k_2+3k_3+\cdots\) số nguyên không vượt quá \(1000\) là bội của ít nhất
một \(a_i\). Các bội này phân biệt, nên
\[
 k_1+2k_2+3k_3+\cdots<1000.
\]
Suy ra
\[
 2k_1+3k_2+4k_3+\cdots
 <1000+n<2000.
\]
Vì thế
\[
 \sum_{i=1}^n\frac1{a_i}
 \le k_1\frac2{1000}+k_2\frac3{1000}+k_3\frac4{1000}+\cdots<2.
\]
Ghi chú: cũng có thể chứng minh \(n\le500\) bằng cách xét ước lẻ lớn nhất;
từ đó còn suy ra \(\sum 1/a_i<3/2\).

\paragraph{Sự kiện hữu ích.}
Định lý Dirichlet: nếu \(a,b\) là các số nguyên dương nguyên tố cùng nhau, thì
cấp số cộng \(a,a+b,a+2b,\ldots\) chứa vô hạn số nguyên tố.

\paragraph{Bài tập.}
\begin{enumerate}
\item Ký hiệu \((a,b,\ldots,g)\) và \([a,b,\ldots,g]\) lần lượt là ước chung
lớn nhất và bội chung nhỏ nhất của các số nguyên dương \(a,b,\ldots,g\).
Chứng minh
\[
 \frac{[a,b,c]^2}{[a,b][a,c][b,c]}
 =
 \frac{(a,b,c)^2}{(a,b)(a,c)(b,c)}.
\]
(USAMO 1972)
\item Chứng minh rằng
\[
 \gcd(a^m-1,a^n-1)=a^{\gcd(m,n)}-1
\]
với mọi số nguyên dương \(a>1,m,n\).
\item Cho \(a,b,c\) là các số nguyên dương. Chứng minh
\[
 \lcm(a,b,c)=\frac{abc\,\gcd(a,b,c)}
 {\gcd(a,b)\gcd(a,c)\gcd(b,c)}.
\]
Biểu diễn \(\gcd(a,b,c)\) theo \(abc,\lcm(a,b,c),\lcm(a,b),\lcm(a,c)\) và
\(\lcm(b,c)\). Tổng quát hóa.
\item Cho \(a,b\) là các số nguyên dương lẻ. Định nghĩa dãy \((f_n)\) bởi
\(f_1=a\), \(f_2=b\), và với \(n\ge3\), \(f_n\) là ước lẻ lớn nhất của
\(f_{n-1}+f_{n-2}\). Chứng minh rằng \(f_n\) hằng từ một lúc nào đó trở đi và
xác định giá trị cuối cùng theo \(a,b\). (USAMO 1993)
\item Cho \(n\ge a_1>a_2>\cdots>a_k\) là các số nguyên dương sao cho
\(\lcm(a_i,a_j)\le n\) với mọi \(i,j\). Chứng minh rằng \(ia_i\le n\) với
\(i=1,2,\ldots,k\).
\end{enumerate}

\section{Các hàm số học}

Có vài hàm số học quan trọng; ở đây ta trình bày ba hàm. Nếu phân tích nguyên
tố của \(n>1\) là \(p_1^{e_1}p_2^{e_2}\cdots p_k^{e_k}\), thì số các số
nguyên dương nhỏ hơn \(n\) và nguyên tố cùng nhau với \(n\) là
\[
\begin{aligned}
\phi(n)
&=\left(1-\frac1{p_1}\right)\left(1-\frac1{p_2}\right)\cdots
\left(1-\frac1{p_k}\right)n\\
&=p_1^{e_1-1}p_2^{e_2-1}\cdots p_k^{e_k-1}
(p_1-1)(p_2-1)\cdots(p_k-1),
\end{aligned}
\]
số ước dương của \(n\) là
\[
 \tau(n)=(e_1+1)(e_2+1)\cdots(e_k+1),
\]
và tổng các ước dương của \(n\) là
\[
\begin{aligned}
\sigma(n)
&=(p_1^{e_1}+p_1^{e_1-1}+\cdots+1)\cdots
(p_k^{e_k}+p_k^{e_k-1}+\cdots+1)\\
&=\frac{p_1^{e_1+1}-1}{p_1-1}\frac{p_2^{e_2+1}-1}{p_2-1}
\cdots\frac{p_k^{e_k+1}-1}{p_k-1}.
\end{aligned}
\]
Ta đặt \(\phi(1)=\tau(1)=\sigma(1)=1\). Một hàm \(f\) được gọi là nhân tính nếu
\(f(mn)=f(m)f(n)\) với mọi cặp số nguyên dương nguyên tố cùng nhau \(m,n\), và
\(f(1)=1\).

\paragraph{Định lý 3.1.}
Các hàm \(\phi,\tau,\sigma\) đều nhân tính. Do đó các công thức trên thu được
dễ dàng bằng cách xét phân tích nguyên tố và tính tại từng lũy thừa nguyên tố.

\paragraph{Ví dụ 3.1.}
Tìm số nghiệm theo cặp có thứ tự số nguyên dương \((x,y)\) của phương trình
\[
 \frac1x+\frac1y=\frac1n,
\]
trong đó \(n\) là số nguyên dương.

\paragraph{Lời giải.}
Ta có
\[
 \frac1x+\frac1y=\frac1n
 \Longleftrightarrow xy=nx+ny
 \Longleftrightarrow (x-n)(y-n)=n^2.
\]
Nếu \(n=1\), nghiệm duy nhất là \((2,2)\). Với
\(n=p_1^{e_1}\cdots p_k^{e_k}\), vì \(x,y>n\), các nghiệm \((x,y)\) tương ứng
song ánh với các ước dương của \(n^2\). Vì vậy số nghiệm là
\[
 \tau(n^2)=(2e_1+1)(2e_2+1)\cdots(2e_k+1).
\]

\paragraph{Ví dụ 3.2.}
Với số nguyên dương \(n\), chứng minh
\[
 \sum_{d\mid n}\phi(d)=n.
\]

\paragraph{Lời giải.}
Với một ước \(d\) của \(n\), gọi \(S_d\) là tập các \(a\), \(1\le a\le n\), sao
cho \(\gcd(a,n)=n/d\). Khi đó \(S_d\) gồm các phần tử dạng \(b\cdot n/d\),
trong đó \(0\le b\le d\) và \(\gcd(b,d)=1\), nên có \(\phi(d)\) phần tử. Mỗi
số từ \(1\) đến \(n\) thuộc đúng một \(S_d\), nên cộng theo mọi ước \(d\mid n\)
cho ta kết quả.

\paragraph{Bài tập.}
\begin{enumerate}
\item Với số nguyên dương \(n\), chứng minh
\[
 \sum_{k=1}^n\tau(k)=\sum_{k=1}^n\left\lfloor\frac nk\right\rfloor.
\]
\item Với số nguyên dương \(n\), chứng minh
\[
 \sum_{d\mid n}\tau^3(d)=\left(\sum_{d\mid n}\tau(d)\right)^2.
\]
\item Chứng minh rằng nếu \(\sigma(N)=2N+1\), thì \(N\) là bình phương của một
số nguyên lẻ. (Putnam 1976)
\end{enumerate}

\section{Số học modulo}

Với số nguyên dương \(m\) và các số nguyên \(a,b\), ta nói \(a\) đồng dư với
\(b\) modulo \(m\) nếu \(m\mid(a-b)\), ký hiệu \(a\equiv b\pmod m\). Nếu
không, ta viết \(a\not\equiv b\pmod m\). Số \(m\) gọi là môđun.

\paragraph{Định lý 4.1.}
Nếu \(a\equiv b\) và \(c\equiv d\pmod m\), thì
\[
 a+c\equiv b+d\pmod m,\qquad ac\equiv bd\pmod m.
\]

\paragraph{Chứng minh.}
Viết \(a=b+km\), \(c=d+\ell m\). Khi đó
\[
 a+c=b+d+(k+\ell)m
\]
và
\[
 ac=bd+(dk+b\ell+k\ell m)m,
\]
nên hai đồng dư thức suy ra.

\paragraph{Các sự kiện hữu ích.}
\begin{itemize}
\item Với mọi số nguyên \(n\),
\[
 n^2\equiv
 \begin{cases}
 0\pmod4,& n\text{ chẵn},\\
 1\pmod4,& n\text{ lẻ}.
 \end{cases}
\]
\item Với mọi số nguyên \(n\),
\[
 n^2\equiv
 \begin{cases}
 0\pmod8,& n\equiv0\pmod4,\\
 4\pmod8,& n\equiv2\pmod4,\\
 1\pmod8,& n\equiv1\pmod2.
 \end{cases}
\]
\item Nếu \(f\) là đa thức hệ số nguyên và \(a\equiv b\pmod m\), thì
\(f(a)\equiv f(b)\pmod m\).
\item Nếu \(f\) là đa thức hệ số nguyên bậc \(n\), không đồng nhất bằng \(0\),
và \(p\) là số nguyên tố, thì đồng dư \(f(x)\equiv0\pmod p\) có nhiều nhất
\(n\) nghiệm modulo \(p\), tính cả bội.
\end{itemize}

\paragraph{Ví dụ 4.1.}
Chứng minh rằng nghiệm hữu tỉ duy nhất của
\[
 x^3+3y^3+9z^3-9xyz=0
\]
là \(x=y=z=0\). (Kürschák 1983)

\paragraph{Lời giải.}
Giả sử có nghiệm hữu tỉ với ít nhất một biến khác \(0\). Vì phương trình thuần
nhất, nhân với lập phương của bội chung nhỏ nhất các mẫu cho ta nghiệm nguyên
\((x_0,y_0,z_0)\). Lấy modulo \(3\), ta được \(x_0^3\equiv0\pmod3\), nên
\(x_0=3x_1\). Thế vào,
\[
 27x_1^3+3y_0^3+9z_0^3-27x_1y_0z_0=0,
\]
suy ra
\[
 y_0^3+3z_0^3+9x_1^3-9x_1y_0z_0=0.
\]
Vậy \((y_0,z_0,x_1)\) lại là nghiệm. Lặp lại, ta buộc \(y_0=3y_1\),
\(z_0=3z_1\), rồi nhận nghiệm \((x_1,y_1,z_1)\). Do đó nghiệm nguyên ban đầu
có thể chia cho lũy thừa \(3\) tùy ý, mâu thuẫn.

\paragraph{Ví dụ 4.2.}
Trong \(10^8+1\) số Fibonacci đầu tiên, có số nào tận cùng bằng \(4\) chữ số
\(0\) không?

\paragraph{Lời giải.}
Có. Xét dãy cặp \((F_k,F_{k+1})\) modulo \(10^4\). Chỉ có hữu hạn khả năng
cho các cặp này, đúng \(10^8\) khả năng, và mỗi cặp phụ thuộc vào cặp trước.
Dãy vì thế tuần hoàn từ một lúc nào đó. Hơn nữa, từ quan hệ Fibonacci ta truy
ngược được cặp trước, nên dãy tuần hoàn ngay từ đầu. Vì \(F_0\equiv0\pmod{10^4}\),
trong vòng \(10^8\) bước phải xuất hiện một số Fibonacci khác chia hết cho
\(10^4\). Thực tế, \(10^4\mid F_{7500}\) và chu kỳ modulo \(10^4\) là \(15000\).

Nếu \(ax\equiv1\pmod m\), ta gọi \(x\) là nghịch đảo của \(a\) modulo \(m\),
ký hiệu \(a^{-1}\); nó là duy nhất modulo \(m\).

\paragraph{Định lý 4.2.}
Nghịch đảo của \(a\) modulo \(m\) tồn tại và duy nhất khi và chỉ khi \(a\) và
\(m\) nguyên tố cùng nhau.

\paragraph{Chứng minh.}
Nếu \(ax\equiv1\pmod m\), thì \(ax=1+km\), nên \(ax-km=1\). Theo Hệ quả 2.2,
\(a,m\) nguyên tố cùng nhau. Ngược lại, nếu \(\gcd(a,m)=1\), tồn tại \(x,y\)
sao cho \(ax+my=1\), nên \(ax\equiv1\pmod m\). Nếu \(x'\) cũng là nghịch đảo,
thì nhân \(ax\equiv ax'\equiv1\) với \(x\) cho \(x\equiv x'\pmod m\).

\paragraph{Hệ quả 4.3.}
Nếu \(p\) là số nguyên tố, thì nghịch đảo của \(a\) modulo \(p\) tồn tại và
duy nhất khi và chỉ khi \(p\nmid a\).

\paragraph{Hệ quả 4.4.}
Nếu \(ak\equiv bk\pmod m\) và \(k\) nguyên tố cùng nhau với \(m\), thì
\(a\equiv b\pmod m\).

Một tập \(\{a_1,\ldots,a_m\}\) là hệ thặng dư đầy đủ modulo \(m\) nếu với mọi
\(i\), \(0\le i\le m-1\), tồn tại duy nhất \(j\) sao cho \(a_j\equiv i\pmod m\).

\paragraph{Ví dụ 4.3.}
Tìm mọi số nguyên dương \(n\) sao cho tồn tại hai hệ thặng dư đầy đủ
\(\{a_i\}\), \(\{b_i\}\) modulo \(n\) mà \(\{a_1+b_1,\ldots,a_n+b_n\}\) cũng
là hệ thặng dư đầy đủ.

\paragraph{Lời giải.}
Đáp án là mọi \(n\) lẻ. Với mọi hệ thặng dư đầy đủ \(\{a_i\}\),
\[
 a_1+\cdots+a_n\equiv\frac{n(n+1)}2\pmod n.
\]
Nếu cả ba tập đều là hệ thặng dư đầy đủ, thì tổng hai hệ đầu đồng dư
\(n(n+1)\equiv0\pmod n\), còn tổng hệ thứ ba đồng dư \(n(n+1)/2\pmod n\).
Vì vậy \(n(n+1)/2\equiv0\pmod n\), điều này buộc \(n\) lẻ. Ngược lại, nếu
\(n\) lẻ, lấy \(a_i=b_i=i\). Khi đó \(a_i+b_i=2i\), và \(2\) nguyên tố cùng
nhau với \(n\), nên \(\{2,4,\ldots,2n\}\) là hệ thặng dư đầy đủ.

\paragraph{Định lý 4.5 (Euler).}
Nếu \(a\) nguyên tố cùng nhau với \(m\), thì
\[
 a^{\phi(m)}\equiv1\pmod m.
\]

\paragraph{Chứng minh.}
Gọi \(a_1,\ldots,a_{\phi(m)}\) là các số nguyên dương nhỏ hơn \(m\) và nguyên
tố cùng nhau với \(m\). Các số \(aa_1,\ldots,aa_{\phi(m)}\) tạo thành một hoán
vị của các \(a_i\) modulo \(m\): chúng đều nguyên tố cùng nhau với \(m\), và
\(aa_i\equiv aa_j\) suy ra \(a_i\equiv a_j\). Do đó
\[
 a_1a_2\cdots a_{\phi(m)}
 \equiv a^{\phi(m)}a_1a_2\cdots a_{\phi(m)}\pmod m,
\]
và sau khi khử tích bên trái bằng nghịch đảo, ta được kết quả.

Định lý này cho công thức \(a^{-1}\equiv a^{\phi(m)-1}\pmod m\); cũng có thể
dùng thuật toán Euclid để tìm nghịch đảo.

\paragraph{Hệ quả 4.6 (định lý nhỏ Fermat).}
Nếu \(p\) là số nguyên tố và \(p\nmid a\), thì \(a^{p-1}\equiv1\pmod p\).

\paragraph{Ví dụ 4.4.}
Chứng minh rằng nếu \(a,b\) là các số nguyên dương nguyên tố cùng nhau, thì tồn
tại các số nguyên \(m,n\) sao cho
\[
 a^m+b^n\equiv1\pmod{ab}.
\]

\paragraph{Lời giải.}
Lấy \(m=\phi(b)\), \(n=\phi(a)\), và đặt \(S=a^m+b^n\). Theo định lý Euler,
\(S\equiv b^{\phi(a)}\equiv1\pmod a\) và \(S\equiv a^{\phi(b)}\equiv1\pmod b\).
Vì \(a,b\) nguyên tố cùng nhau, \(S\equiv1\pmod{ab}\).

\paragraph{Ví dụ 4.5.}
Với số nguyên dương \(i\), gọi \(S_i\) là tổng các tích chập \(i\) phần tử của
\(1,2,\ldots,p-1\), trong đó \(p\) là số nguyên tố lẻ. Chứng minh
\[
 S_1\equiv S_2\equiv\cdots\equiv S_{p-2}\equiv0\pmod p.
\]

\paragraph{Lời giải.}
Ta có
\[
 (x-1)(x-2)\cdots(x-(p-1))
 =x^{p-1}-S_1x^{p-2}+S_2x^{p-3}-\cdots-S_{p-2}x+S_{p-1}.
\]
Đa thức này triệt tiêu tại \(x=1,\ldots,p-1\). Theo định lý nhỏ Fermat,
\(x^{p-1}-1\) cũng vậy modulo \(p\). Hiệu của hai đa thức có bậc \(p-2\) nhưng
có \(p-1\) nghiệm modulo \(p\), nên phải là đa thức không; so sánh hệ số cho
kết quả. Đồng thời \(S_{p-1}=(p-1)!\equiv-1\pmod p\), chính là định lý Wilson.

\paragraph{Định lý 4.7.}
Nếu \(p\) là số nguyên tố và \(p\mid(4n^2+1)\), thì \(p\equiv1\pmod4\).

\paragraph{Chứng minh.}
Rõ ràng \(p\ne2\). Giả sử \(p=4k+3\). Đặt \(y=2n\). Khi đó \(y^2+1\equiv0\)
và theo Fermat,
\[
 y^{p-1}=y^{4k+2}=(y^2)^{2k+1}\equiv(-1)^{2k+1}\equiv-1\pmod p,
\]
mâu thuẫn với \(y^{p-1}\equiv1\pmod p\). Vậy \(p\equiv1\pmod4\).

\paragraph{Ví dụ 4.6.}
Chứng minh rằng có vô hạn số nguyên tố dạng \(4k+1\) và vô hạn số nguyên tố
dạng \(4k+3\).

\paragraph{Lời giải.}
Nếu chỉ có hữu hạn số nguyên tố dạng \(4k+1\), là \(p_1,\ldots,p_n\), đặt
\[
 N=4(p_1p_2\cdots p_n)^2+1.
\]
Theo Định lý 4.7, mọi ước nguyên tố của \(N\) có dạng \(4k+1\), nhưng \(N\)
không chia hết cho bất kỳ \(p_i\) nào, mâu thuẫn. Tương tự, nếu chỉ có hữu hạn
số nguyên tố dạng \(4k+3\), là \(q_1,\ldots,q_m\), đặt
\[
 M=4q_1q_2\cdots q_m-1.
\]
Khi đó \(M\equiv3\pmod4\), nên \(M\) có một ước nguyên tố dạng \(4k+3\), nhưng
không chia hết cho bất kỳ \(q_i\) nào, mâu thuẫn.

\paragraph{Ví dụ 4.7.}
Chứng minh rằng nếu \(n>1\), thì \(n\nmid(2^n-1)\).

\paragraph{Lời giải.}
Gọi \(p\) là ước nguyên tố nhỏ nhất của \(n\). Khi đó \(\gcd(n,p-1)=1\), nên
tồn tại \(x,y\) sao cho \(nx+(p-1)y=1\). Nếu \(p\mid(2^n-1)\), thì
\[
 2\equiv2^{nx+(p-1)y}\equiv(2^n)^x(2^{p-1})^y\equiv1\pmod p,
\]
mâu thuẫn. Vậy \(p\nmid(2^n-1)\), nên \(n\nmid(2^n-1)\).

\paragraph{Định lý 4.8 (Wilson).}
Nếu \(p\) là số nguyên tố, thì
\[
 (p-1)!\equiv-1\pmod p.
\]

\paragraph{Chứng minh.}
Đồng dư \(x^2\equiv1\pmod p\) chỉ có nghiệm \(x\equiv\pm1\). Vì vậy mỗi
\(i=2,\ldots,p-2\) ghép cặp với nghịch đảo riêng biệt của nó trong cùng khoảng.
Do đó
\[
 (p-1)!\equiv1\cdot1\cdots1\cdot(p-1)\equiv-1\pmod p.
\]

\paragraph{Ví dụ 4.8.}
Cho hai hệ thặng dư đầy đủ modulo \(101\):
\[
 \{a_1,\ldots,a_{101}\},\qquad \{b_1,\ldots,b_{101}\}.
\]
Tập
\[
 \{a_1b_1,\ldots,a_{101}b_{101}\}
\]
có thể là hệ thặng dư đầy đủ modulo \(101\) không?

\paragraph{Lời giải.}
Không. Giả sử có thể. Không mất tính tổng quát, \(a_{101}\equiv0\pmod{101}\).
Khi đó \(b_{101}\equiv0\pmod{101}\), nếu không hai tích sẽ cùng đồng dư \(0\).
Theo Wilson,
\[
 a_1\cdots a_{100}\equiv b_1\cdots b_{100}\equiv100!\equiv-1\pmod{101},
\]
nên \(a_1b_1\cdots a_{100}b_{100}\equiv1\pmod{101}\). Nhưng nếu các tích cũng
là hệ thặng dư đầy đủ, tích các phần tử khác \(0\) của chúng phải bằng
\(100!\equiv-1\), mâu thuẫn.

\paragraph{Định lý 4.9.}
Nếu \(p\) là số nguyên tố, thì \(x^2+1\equiv0\pmod p\) có nghiệm khi và chỉ
khi \(p=2\) hoặc \(p\equiv1\pmod4\).

\paragraph{Chứng minh.}
Với \(p=2\), \(x=1\) là nghiệm. Với \(p\equiv3\pmod4\), không có nghiệm theo
nhận xét từ Định lý 4.7. Với \(p=4k+1\), đặt \(x=1\cdot2\cdots(2k)\). Khi đó
\[
\begin{aligned}
x^2&\equiv1\cdot2\cdots(2k)\cdot(2k)\cdots2\cdot1\\
&\equiv1\cdot2\cdots(2k)\cdot(-2k)\cdots(-2)(-1)\\
&\equiv(p-1)!\equiv-1\pmod p.
\end{aligned}
\]

\paragraph{Định lý 4.10.}
Nếu \(p\) là số nguyên tố và \(p\equiv1\pmod4\), thì tồn tại các số nguyên
dương \(x,y\) sao cho \(p=x^2+y^2\).

\paragraph{Chứng minh.}
Theo Định lý 4.9, tồn tại \(a\) sao cho \(a^2\equiv-1\pmod p\). Xét các số
\(ax-y\), với \(0\le x,y<\sqrt p\). Có \((\lfloor\sqrt p\rfloor+1)^2>p\) cặp,
nên tồn tại hai cặp khác nhau cho cùng phần dư. Lấy hiệu được \(x,y\), không
đồng thời bằng \(0\), sao cho \(ax\equiv y\pmod p\). Do đó
\[
 x^2+y^2\equiv x^2+a^2x^2\equiv0\pmod p.
\]
Mặt khác \(0<x^2+y^2<2p\), nên \(x^2+y^2=p\).

\paragraph{Định lý 4.11.}
Cho \(n\) là số nguyên dương. Tồn tại các số nguyên \(x,y\) sao cho
\(n=x^2+y^2\) khi và chỉ khi mỗi ước nguyên tố của \(n\) dạng \(4k+3\) xuất
hiện với số mũ chẵn.

\paragraph{Định lý 4.12 (định lý phần dư Trung Hoa).}
Nếu \(a_1,\ldots,a_k\) là các số nguyên, và \(m_1,\ldots,m_k\) là các số
nguyên đôi một nguyên tố cùng nhau, thì hệ
\[
 x\equiv a_1\pmod{m_1},\quad
 x\equiv a_2\pmod{m_2},\quad\ldots,\quad
 x\equiv a_k\pmod{m_k}
\]
có nghiệm duy nhất modulo \(m_1m_2\cdots m_k\).

\paragraph{Chứng minh.}
Đặt \(m=m_1\cdots m_k\). Vì \(m/m_i\) nguyên tố cùng nhau với \(m_i\), tồn tại
\(t_i\) sao cho \(t_i(m/m_i)\equiv1\pmod{m_i}\). Đặt
\(s_i=t_i(m/m_i)\). Khi đó \(s_i\equiv1\pmod{m_i}\) và
\(s_i\equiv0\pmod{m_j}\) nếu \(j\ne i\). Vậy
\[
 x=a_1s_1+\cdots+a_ks_k
\]
là nghiệm. Nếu \(x'\) là nghiệm khác, thì \(x-x'\equiv0\pmod{m_i}\) với mọi
\(i\), nên \(x\equiv x'\pmod m\).

\paragraph{Ví dụ 4.9.}
Với số nguyên dương \(n\), tìm số nghiệm của \(x^2\equiv1\pmod n\).

\paragraph{Lời giải.}
Viết \(n=2^e p_1^{e_1}\cdots p_k^{e_k}\). Theo định lý phần dư Trung Hoa, số
nghiệm là tích số nghiệm theo từng lũy thừa nguyên tố. Với \(p_i\) lẻ,
\[
 x^2\equiv1\pmod{p_i^{e_i}}
\]
tương đương \((x-1)(x+1)\equiv0\), và \(p_i\) không thể chia cả hai thừa số,
nên có hai nghiệm \(x\equiv\pm1\pmod{p_i^{e_i}}\). Với \(2^e\), nếu \(e=1,2\)
có lần lượt \(1,2\) nghiệm. Nếu \(e\ge3\), các nghiệm là
\[
 x\equiv1,\;2^{e-1}-1,\;2^{e-1}+1,\;2^e-1\pmod{2^e}.
\]
Vậy số nghiệm là
\[
\begin{array}{c|c}
e&\text{số nghiệm}\\ \hline
0,1&2^k\\
2&2^{k+1}\\
\ge3&2^{k+2}.
\end{array}
\]

\paragraph{Định lý 4.13.}
Cho \(m\) là số nguyên dương, \(a,b\) là các số nguyên, và
\(k=\gcd(a,m)\). Đồng dư \(ax\equiv b\pmod m\) có \(k\) nghiệm hoặc không có
nghiệm tùy theo \(k\mid b\) hoặc \(k\nmid b\).

\paragraph{Bài tập.}
\begin{enumerate}
\item Chứng minh rằng với mỗi số nguyên dương \(n\), tồn tại \(n\) số nguyên
dương liên tiếp sao cho không số nào là lũy thừa nguyên của một số nguyên tố.
(IMO 1989)
\item Với số nguyên dương lẻ \(n>1\), gọi \(S\) là tập các số nguyên \(x\),
\(1\le x\le n\), sao cho cả \(x\) và \(x+1\) đều nguyên tố cùng nhau với \(n\).
Chứng minh
\[
 \prod_{x\in S}x\equiv1\pmod n.
\]
\item Tìm tất cả nghiệm nguyên dương của \(3^x+4^y=5^z\). (Danh sách chọn lọc
IMO 1991)
\item Cho \(n\) là số nguyên dương sao cho \(n+1\) chia hết cho \(24\). Chứng
minh tổng các ước dương của \(n\) chia hết cho \(24\). (Putnam 1969)
\item (Định lý Wolstenholme) Chứng minh rằng nếu
\[
 1+\frac12+\frac13+\cdots+\frac1{p-1}
\]
được viết thành phân số, với \(p\ge5\) nguyên tố, thì \(p^2\) chia tử số.
\item Cho \(a\) là nghiệm dương lớn nhất của \(x^3-3x^2+1=0\). Chứng minh
\(\lfloor a^{1788}\rfloor\) và \(\lfloor a^{1988}\rfloor\) đều chia hết cho
\(17\). (Danh sách chọn lọc IMO 1988)
\item Cho \(\{a_1,\ldots,a_n\}\) và \(\{b_1,\ldots,b_n\}\) là các hệ thặng dư
đầy đủ modulo \(n\), sao cho \(\{a_1b_1,\ldots,a_nb_n\}\) cũng là hệ thặng dư
đầy đủ modulo \(n\). Chứng minh \(n=1\) hoặc \(2\).
\item Cho \(m,n\) là các số nguyên dương. Chứng minh rằng \(4mn-m-n\) không
bao giờ là số chính phương. (Đề xuất IMO 1984)
\end{enumerate}

\section{Hệ số nhị thức}

Với các số nguyên không âm \(n,k\), \(k\le n\), hệ số nhị thức được định nghĩa
bởi
\[
 \binom nk=\frac{n!}{k!(n-k)!},
\]
và quy ước \(\binom nk=0\) nếu \(k>n\). Trong các kết quả sau, với đa thức hệ
số nguyên \(f,g\), ta viết \(f\equiv g\pmod m\) nếu \(m\) chia mọi hệ số của
\(f-g\).

\paragraph{Định lý 5.1.}
Nếu \(p\) là số nguyên tố, thì số thừa số \(p\) trong \(n!\) là
\[
 \left\lfloor\frac np\right\rfloor+
 \left\lfloor\frac n{p^2}\right\rfloor+
 \left\lfloor\frac n{p^3}\right\rfloor+\cdots.
\]
Nó cũng bằng \((n-s_n)/(p-1)\), trong đó \(s_n\) là tổng chữ số của \(n\) khi
viết trong cơ số \(p\).

\paragraph{Định lý 5.2.}
Nếu \(p\) là số nguyên tố, thì
\[
 \binom pi\equiv0\pmod p
\]
với \(1\le i\le p-1\).

\paragraph{Hệ quả 5.3.}
\[
 (1+x)^p\equiv1+x^p\pmod p.
\]

\paragraph{Bổ đề 5.4.}
Với mọi số thực \(x,y\), \(\lfloor x+y\rfloor\ge\lfloor x\rfloor+\lfloor y\rfloor\).

\paragraph{Định lý 5.5.}
Nếu \(p\) là số nguyên tố, thì
\[
 \binom{p^k}{i}\equiv0\pmod p
\]
với \(1\le i\le p^k-1\).

\paragraph{Chứng minh.}
Theo Bổ đề 5.4,
\[
 \sum_{j=1}^k\left(
\left\lfloor\frac i{p^j}\right\rfloor+
\left\lfloor\frac{p^k-i}{p^j}\right\rfloor
\right)
\le
\sum_{j=1}^k\left\lfloor\frac{p^k}{p^j}\right\rfloor.
\]
Hai vế là số thừa số \(p\) trong \(i!(p^k-i)!\) và trong \(p^k!\). Vì
\(\binom{p^k}{i}\) không bằng \(0\) hay \(1\), bất đẳng thức là nghiêm ở ít
nhất một chỗ; do đó \(\binom{p^k}{i}\) chia hết cho \(p\).

\paragraph{Hệ quả 5.6.}
\[
 (1+x)^{p^k}\equiv1+x^{p^k}\pmod p.
\]

\paragraph{Ví dụ 5.1.}
Với số nguyên dương \(n\), chứng minh tích của \(n\) số nguyên dương liên tiếp
chia hết cho \(n!\).

\paragraph{Lời giải.}
Nếu các số là \(m,m+1,\ldots,m+n-1\), thì
\[
 \frac{m(m+1)\cdots(m+n-1)}{n!}=\binom{m+n-1}{n}.
\]

\paragraph{Ví dụ 5.2.}
Với số nguyên dương \(n\), chứng minh
\[
 (n+1)\lcm\left(\binom n0,\binom n1,\ldots,\binom nn\right)
 =\lcm(1,2,\ldots,n+1).
\]

\paragraph{Lời giải.}
Cho \(p\le n+1\) là số nguyên tố. Gọi \(\alpha,\beta\) lần lượt là số mũ cao
nhất của \(p\) trong vế trái và vế phải. Chọn \(r\) sao cho
\(p^r\le n+1<p^{r+1}\); rõ ràng \(\beta=r\). Ta dùng nhận xét:
\[
 \text{nếu }p^r\le m<p^{r+1},\text{ thì }p^{r+1}\nmid\binom mk
 \quad(0\le k\le m).
\]
Thật vậy, số mũ của \(p\) trong \(\binom mk\) là
\[
 \gamma=\sum_{s=1}^r\left(
\left\lfloor\frac m{p^s}\right\rfloor-
\left\lfloor\frac k{p^s}\right\rfloor-
\left\lfloor\frac{m-k}{p^s}\right\rfloor\right),
\]
mỗi hạng là \(0\) hoặc \(1\), nên \(\gamma\le r\).

Đặt
\[
 a_k=(n+1)\binom nk=(n-k+1)\binom{n+1}k=(k+1)\binom{n+1}{k+1}.
\]
Theo nhận xét trên, \(p^{r+1}\) không chia bất kỳ \(\binom nk\),
\(\binom{n+1}k\), \(\binom{n+1}{k+1}\). Vì vậy \(p^{r+1}\) chỉ có thể chia
\(a_k\) nếu \(p\) chia đồng thời \(n+1,n-k+1,k+1\), điều này cho
\(p\mid-1\), vô lý. Mặt khác, với \(k=p^r-1\), \(a_k=(k+1)\binom{n+1}{k+1}\)
chia hết cho \(p^r\). Vậy \(\alpha=\beta=r\).

\paragraph{Định lý 5.7 (Lucas).}
Cho \(m,n\) là các số nguyên không âm và \(p\) là số nguyên tố. Viết
\[
 m=m_kp^k+m_{k-1}p^{k-1}+\cdots+m_1p+m_0,
\]
\[
 n=n_kp^k+n_{k-1}p^{k-1}+\cdots+n_1p+n_0
\]
trong cơ số \(p\). Khi đó
\[
 \binom mn\equiv
 \binom{m_k}{n_k}\binom{m_{k-1}}{n_{k-1}}\cdots
 \binom{m_1}{n_1}\binom{m_0}{n_0}\pmod p.
\]

\paragraph{Chứng minh.}
Theo Hệ quả 5.6,
\[
\begin{aligned}
(1+x)^m
&\equiv(1+x^{p^k})^{m_k}(1+x^{p^{k-1}})^{m_{k-1}}\cdots
(1+x^p)^{m_1}(1+x)^{m_0}\pmod p.
\end{aligned}
\]
So sánh hệ số của \(x^n\) trên hai vế cho kết quả.

\paragraph{Hệ quả 5.8.}
Cho \(n\) là số nguyên dương. Gọi \(A(n)\) là số thừa số \(2\) trong \(n!\),
và \(B(n)\) là số chữ số \(1\) trong biểu diễn nhị phân của \(n\). Số hệ số lẻ
ở hàng thứ \(n\) của tam giác Pascal, tương đương số hệ số lẻ trong khai triển
\((1+x)^n\), là \(2^{B(n)}\). Hơn nữa, \(A(n)+B(n)=n\).

\paragraph{Sự kiện hữu ích.}
Với đa thức hệ số nguyên \(f\) và số nguyên tố \(p\),
\[
 [f(x)]^{p^n}\equiv f(x^{p^n})\pmod p.
\]

\paragraph{Bài tập.}
\begin{enumerate}
\item Cho \(a,b\) là các số nguyên không âm và \(p\) là số nguyên tố. Chứng
minh
\[
 \binom{pa}{pb}\equiv\binom ab\pmod p.
\]
\item Gọi \(a_n\) là chữ số khác \(0\) cuối cùng trong biểu diễn thập phân của
\(n!\). Dãy \(a_1,a_2,a_3,\ldots\) có tuần hoàn từ một lúc nào đó không?
(Danh sách chọn lọc IMO 1991)
\item Tìm mọi số nguyên dương \(n\) sao cho \(2^n\mid(3^n-1)\).
\item Tìm số nguyên lớn nhất \(k\) sao cho \(1991^k\) chia
\[
 1990^{1991^{1992}}+1992^{1991^{1990}}.
\]
(Danh sách chọn lọc IMO 1991)
\item Với số nguyên dương \(n\), gọi \(a(n)\) và \(b(n)\) là số hệ số nhị thức
ở hàng thứ \(n\) của tam giác Pascal lần lượt đồng dư \(1\) và \(2\) modulo
\(3\). Chứng minh \(a(n)-b(n)\) luôn là một lũy thừa của \(2\).
\item Cho \(n\) là số nguyên dương. Chứng minh rằng nếu số thừa số \(2\) trong
\(n!\) là \(n-1\), thì \(n\) là lũy thừa của \(2\).
\item Với số nguyên dương \(n\), đặt
\[
 C_n=\frac1{n+1}\binom{2n}{n},\qquad S_n=C_1+C_2+\cdots+C_n.
\]
Chứng minh \(S_n\equiv1\pmod3\) khi và chỉ khi trong biểu diễn cơ số \(3\) của
\(n+1\) có một chữ số \(2\).
\end{enumerate}

\section{Cấp của một phần tử}

Nếu \(a\) nguyên tố cùng nhau với \(m\), tồn tại số nguyên dương \(n\) sao cho
\(a^n\equiv1\pmod m\). Gọi \(d\) là số nhỏ nhất như vậy. Ta gọi \(d\) là cấp
của \(a\) modulo \(m\), ký hiệu \(\ord_m(a)\), hoặc \(\ord(a)\) nếu môđun đã
rõ.

\paragraph{Định lý 6.1.}
Nếu \(a\) nguyên tố cùng nhau với \(m\), thì \(a^n\equiv1\pmod m\) khi và chỉ
khi \(\ord(a)\mid n\). Hơn nữa,
\[
 a^{n_0}\equiv a^{n_1}\pmod m
 \quad\Longleftrightarrow\quad
 \ord(a)\mid(n_0-n_1).
\]

\paragraph{Chứng minh.}
Đặt \(d=\ord(a)\). Nếu \(d\mid n\) thì rõ. Ngược lại, viết \(n=qd+r\),
\(0\le r<d\). Khi đó \(a^n\equiv(a^d)^qa^r\equiv a^r\equiv1\pmod m\). Do tính
nhỏ nhất của \(d\), \(r=0\). Phần thứ hai suy ra bằng cách nhân với nghịch đảo.
Đặc biệt, theo định lý Euler, \(\ord(a)\mid\phi(m)\).

\paragraph{Ví dụ 6.1.}
Chứng minh cấp của \(2\) modulo \(101\) là \(100\).

\paragraph{Lời giải.}
Gọi \(d=\ord(2)\). Khi đó \(d\mid100\). Nếu \(d<100\), thì \(d\) chia \(50\)
hoặc \(20\). Nhưng
\[
 2^{50}\equiv1024^5\equiv14^5\equiv196\cdot196\cdot14
\equiv(-6)(-6)14\equiv-1\pmod{101},
\]
và
\[
 2^{20}\equiv1024^2\equiv14^2\equiv-6\pmod{101}.
\]
Vậy \(d=100\).

\paragraph{Ví dụ 6.2.}
Chứng minh rằng nếu \(p\) là số nguyên tố, thì mọi ước nguyên tố của \(2^p-1\)
đều lớn hơn \(p\).

\paragraph{Lời giải.}
Cho \(q\mid(2^p-1)\) nguyên tố. Khi đó \(2^p\equiv1\pmod q\), nên
\(\ord(2)\mid p\). Vì \(\ord(2)\ne1\), ta có \(\ord(2)=p\). Theo định lý nhỏ
Fermat, \(\ord(2)\mid(q-1)\), nên \(p\le q-1\). Thực ra với \(p>2\), \(q\) có
dạng \(2kp+1\).

\paragraph{Ví dụ 6.3.}
Cho \(p\) là số nguyên tố nguyên tố cùng nhau với \(10\), và \(0<n<p\). Gọi
\(d\) là cấp của \(10\) modulo \(p\).
\begin{enumerate}[label=(\alph*)]
\item Chứng minh độ dài chu kỳ thập phân của \(n/p\) là \(d\).
\item Nếu \(d\) chẵn, chứng minh chu kỳ thập phân của \(n/p\) có thể chia
thành hai nửa có tổng bằng \(10^{d/2}-1\). Ví dụ
\(1/7=0.\overline{142857}\), \(142+857=999=10^3-1\).
\end{enumerate}

\paragraph{Lời giải.}
(a) Gọi \(m\) là độ dài chu kỳ và viết
\[
 \frac np=0.\overline{a_1a_2\cdots a_m}.
\]
Khi đó
\[
 \frac{(10^m-1)n}{p}=a_1a_2\cdots a_m
\]
là số nguyên, nên \(p\mid10^m-1\), do đó \(d\mid m\). Ngược lại,
\(p\mid10^d-1\), nên \((10^d-1)n/p\) là số nguyên có không quá \(d\) chữ số;
suy ra chu kỳ có độ dài không quá \(d\). Vậy \(m=d\).

(b) Viết \(d=2k\). Vì \(p\mid10^{2k}-1=(10^k-1)(10^k+1)\) nhưng
\(p\nmid10^k-1\), ta có \(p\mid10^k+1\). Từ
\[
 \frac{10^kn}{p}=a_1\cdots a_k.a_{k+1}\cdots a_{2k}
\]
suy ra
\[
 \frac{(10^k+1)n}{p}
 =a_1\cdots a_k+0.a_1\cdots a_k+0.a_{k+1}\cdots a_{2k}
\]
là số nguyên. Điều này xảy ra khi và chỉ khi
\[
 a_1\cdots a_k+a_{k+1}\cdots a_{2k}=10^k-1.
\]

\paragraph{Bài tập.}
\begin{enumerate}
\item Chứng minh rằng với mọi số nguyên dương \(a>1\) và \(n\),
\[
 n\mid\phi(a^n-1).
\]
\item Chứng minh rằng nếu \(p\) là số nguyên tố, thì \(p^p-1\) có một ước
nguyên tố đồng dư \(1\) modulo \(p\).
\item Với mọi số nguyên \(a\), đặt \(n_a=101^a-100\cdot2^a\). Chứng minh rằng
nếu \(0\le a,b,c,d\le99\) và
\[
 n_a+n_b\equiv n_c+n_d\pmod{10100},
\]
thì \(\{a,b\}=\{c,d\}\). (Putnam 1994)
\item Chứng minh rằng nếu \(3\le d\le2^{n+1}\), thì
\[
 d\nmid(a^{2^n}+1)
\]
với mọi số nguyên dương \(a\).
\end{enumerate}

\section{Thặng dư bậc hai}

Cho \(m>1\) và \(a\) nguyên tố cùng nhau với \(m\). Nếu \(x^2\equiv a\pmod m\)
có nghiệm, ta nói \(a\) là thặng dư bậc hai modulo \(m\); nếu không, \(a\) là
không thặng dư bậc hai. Với \(p\) nguyên tố lẻ, ký hiệu Legendre
\[
 \leg ap
\]
nhận giá trị \(1\) nếu \(a\) là thặng dư bậc hai modulo \(p\), và \(-1\) nếu
không.

\paragraph{Định lý 7.1.}
Cho \(p\) là số nguyên tố lẻ, và \(a,b\) nguyên tố cùng nhau với \(p\). Khi đó
\[
 \leg ap\equiv a^{(p-1)/2}\pmod p,
\]
và
\[
 \leg ap\leg bp=\leg{ab}p.
\]

\paragraph{Chứng minh.}
Nếu \(x^2\equiv a\pmod p\) có nghiệm, thì
\[
 a^{(p-1)/2}\equiv x^{p-1}\equiv1\pmod p.
\]
Nếu không có nghiệm, thì với mỗi \(i=1,\ldots,p-1\), có duy nhất \(j\ne i\)
sao cho \(ij\equiv a\pmod p\). Ghép tất cả các số thành \((p-1)/2\) cặp như
vậy và nhân lại:
\[
 a^{(p-1)/2}\equiv1\cdot2\cdots(p-1)\equiv(p-1)!\equiv-1\pmod p.
\]
Tính nhân tính suy ra từ tiêu chuẩn này. Phần đầu gọi là tiêu chuẩn Euler.

\paragraph{Ví dụ 7.1.}
Chứng minh rằng nếu \(p\) là số nguyên tố lẻ, thì
\[
 \leg1p+\leg2p+\cdots+\leg{p-1}p=0.
\]

\paragraph{Lời giải.}
Các số \(1^2,2^2,\ldots,((p-1)/2)^2\) phân biệt modulo \(p\), còn các bình
phương của \((p+1)/2,\ldots,p-1\) cho cùng các thặng dư theo thứ tự ngược. Vì
thế có đúng \((p-1)/2\) thặng dư và \((p-1)/2\) không thặng dư.

\paragraph{Định lý 7.2 (bổ đề Gauss).}
Cho \(p\) là số nguyên tố lẻ và \(a\) nguyên tố cùng nhau với \(p\). Xét các
thặng dư không âm nhỏ nhất của
\[
 a,2a,\ldots,\frac{p-1}{2}a
\]
modulo \(p\). Nếu \(n\) là số thặng dư trong đó lớn hơn \(p/2\), thì
\[
 \leg ap=(-1)^n.
\]

\paragraph{Định lý 7.3.}
Nếu \(p\) là số nguyên tố lẻ, thì
\[
 \leg{-1}p=(-1)^{(p-1)/2}
 =
\begin{cases}
1,&p\equiv1\pmod4,\\
-1,&p\equiv3\pmod4.
\end{cases}
\]

\paragraph{Định lý 7.4.}
Nếu \(p\) là số nguyên tố lẻ, thì
\[
 \leg2p=(-1)^{(p^2-1)/8}
 =
\begin{cases}
1,&p\equiv1\text{ hoặc }7\pmod8,\\
-1,&p\equiv3\text{ hoặc }5\pmod8.
\end{cases}
\]

\paragraph{Chứng minh.}
Nếu \(p\equiv1\) hoặc \(5\pmod8\), thì
\[
 2^{(p-1)/2}\left(\frac{p-1}{2}\right)!
 \equiv2\cdot4\cdot\ldots\cdot(p-1)
 \equiv(-1)^{(p-1)/4}\left(\frac{p-1}{2}\right)!\pmod p.
\]
Suy ra \(2^{(p-1)/2}\equiv(-1)^{(p-1)/4}\), nên \(\leg2p\) lần lượt là \(1\)
hoặc \(-1\). Trường hợp \(p\equiv3\) hoặc \(7\pmod8\) tương tự, cho
\[
 2^{(p-1)/2}\equiv(-1)^{(p+1)/4}\pmod p,
\]
nên \(\leg2p\) lần lượt là \(-1\) hoặc \(1\).

\paragraph{Ví dụ 7.2.}
Chứng minh rằng nếu \(n\) là số nguyên dương lẻ, thì mọi ước nguyên tố của
\(2^n-1\) có dạng \(8k\pm1\).

\paragraph{Lời giải.}
Cho \(p\mid2^n-1\), \(p\) nguyên tố, và viết \(n=2m+1\). Khi đó
\[
 2^n=2(2^m)^2\equiv1\pmod p,
\]
nên \(2\) là thặng dư bậc hai modulo \(p\). Theo Định lý 7.4, \(p\equiv\pm1
\pmod8\).

\paragraph{Định lý 7.5 (luật thuận nghịch bậc hai).}
Với hai số nguyên tố lẻ phân biệt \(p,q\),
\[
 \leg pq\leg qp=(-1)^{\frac{p-1}{2}\cdot\frac{q-1}{2}}.
\]

\paragraph{Ví dụ 7.3.}
Với những số nguyên tố \(p>3\) nào thì \(x^2\equiv-3\pmod p\) có nghiệm?

\paragraph{Lời giải.}
Ta cần \(\leg{-3}p=1\). Theo thuận nghịch bậc hai,
\[
 \leg3p\leg p3=(-1)^{(p-1)/2}=\leg{-1}p.
\]
Do đó
\[
 \leg{-3}p=\leg{-1}p\leg3p
 =\leg{-1}p\leg p3\leg{-1}p=\leg p3.
\]
Vậy điều kiện là \(p\equiv1\pmod3\). Vì \(p>3\) là nguyên tố, điều này tương
đương \(p\equiv1\pmod6\).

\paragraph{Ví dụ 7.4.}
Chứng minh rằng nếu \(p=2^n+1\), \(n\ge2\), là số nguyên tố, thì
\[
 3^{(p-1)/2}+1
\]
chia hết cho \(p\).

\paragraph{Lời giải.}
Ta phải có \(n\) chẵn, nếu không \(p\equiv0\pmod3\). Theo tiêu chuẩn Euler,
\[
 3^{(p-1)/2}\equiv\leg3p\pmod p.
\]
Vì \(p\equiv1\pmod4\) và \(p\equiv2\pmod3\), nên theo thuận nghịch bậc hai
\[
 \leg3p=\leg p3=-1.
\]
Do đó \(3^{(p-1)/2}+1\equiv0\pmod p\).

\paragraph{Các sự kiện hữu ích.}
\begin{itemize}
\item Nếu \(p\) là số nguyên tố và \(p\equiv1\) hoặc \(3\pmod8\), thì tồn tại
các số nguyên dương \(x,y\) sao cho \(p=x^2+2y^2\).
\item Nếu \(p\) là số nguyên tố và \(p\equiv1\pmod6\), thì tồn tại các số
nguyên dương \(x,y\) sao cho \(p=x^2+3y^2\).
\end{itemize}

\paragraph{Bài tập.}
\begin{enumerate}
\item Chứng minh rằng nếu \(p>3\) là số nguyên tố, thì tổng các thặng dư bậc
hai trong \(1,2,\ldots,p-1\) chia hết cho \(p\).
\item Gọi \(F_n\) là số Fibonacci thứ \(n\). Chứng minh rằng nếu \(p>5\) là số
nguyên tố, thì
\[
 F_p\equiv\leg p5\pmod p.
\]
\item Chứng minh rằng \(16\) là lũy thừa bậc \(8\) hoàn hảo modulo mọi số
nguyên tố \(p\).
\item Cho \(a,b,c\) là các số nguyên dương đôi một nguyên tố cùng nhau và
\(a^2-ab+b^2=c^2\). Chứng minh mọi ước nguyên tố của \(c\) có dạng \(6k+1\).
\item Cho \(p\) là số nguyên tố lẻ và \(\zeta\) là căn nguyên thủy bậc \(p\)
của đơn vị, tức \(\zeta^p=1\) và \(\zeta^k\ne1\) với \(1\le k\le p-1\). Gọi
\(A_p,B_p\) lần lượt là tập các thặng dư và không thặng dư bậc hai modulo \(p\).
Đặt
\[
 \alpha=\sum_{k\in A_p}\zeta^k,\qquad
 \beta=\sum_{k\in B_p}\zeta^k.
\]
Ví dụ, với \(p=7\), \(\alpha=\zeta+\zeta^2+\zeta^4\) và
\(\beta=\zeta^3+\zeta^5+\zeta^6\). Chứng minh \(\alpha,\beta\) là nghiệm của
\[
 x^2+x+\frac{1-\leg{-1}p\,p}{4}=0.
\]
\end{enumerate}

\section{Căn nguyên thủy}

Nếu cấp của \(g\) modulo \(m\) là \(\phi(m)\), ta gọi \(g\) là căn nguyên thủy
modulo \(m\).

\paragraph{Ví dụ 8.1.}
Chứng minh rằng \(2\) là căn nguyên thủy modulo \(3^n\) với mọi \(n\ge1\).

\paragraph{Lời giải.}
Trường hợp \(n=1\) rõ. Giả sử đúng với \(n=k\), tức
\[
 2^{\phi(3^k)}=2^{2\cdot3^{k-1}}\equiv1\pmod{3^k}.
\]
Gọi \(d\) là cấp của \(2\) modulo \(3^{k+1}\). Khi đó \(2\cdot3^{k-1}\mid d\)
và \(d\mid\phi(3^{k+1})=2\cdot3^k\), nên \(d\) là \(2\cdot3^{k-1}\) hoặc
\(2\cdot3^k\). Ta dùng bổ đề:
\[
 2^{2\cdot3^{n-1}}\equiv1+3^n\pmod{3^{n+1}}\quad(n\ge1).
\]
Bổ đề đúng với \(n=1\). Nếu đúng với \(n=k\), viết
\[
 2^{2\cdot3^{k-1}}=1+3^k+3^{k+1}m.
\]
Lập phương hai vế cho
\[
 2^{2\cdot3^k}\equiv1+3^{k+1}\pmod{3^{k+2}}.
\]
Vậy bổ đề đúng. Do đó \(2^{2\cdot3^{k-1}}\not\equiv1\pmod{3^{k+1}}\), nên cấp
là \(2\cdot3^k\).

\paragraph{Hệ quả 8.2.}
Nếu \(2^n\equiv-1\pmod{3^k}\), thì \(3^{k-1}\mid n\).

\paragraph{Định lý 8.3.}
Nếu \(m\) có căn nguyên thủy, thì nó có \(\phi(\phi(m))\) căn nguyên thủy phân
biệt modulo \(m\).

\paragraph{Định lý 8.4.}
Số nguyên dương \(m\) có căn nguyên thủy khi và chỉ khi \(m\) là một trong các
dạng \(2\), \(4\), \(p^k\), hoặc \(2p^k\), trong đó \(p\) là số nguyên tố lẻ.

\paragraph{Định lý 8.5.}
Nếu \(g\) là căn nguyên thủy của \(m\), thì \(g^n\equiv1\pmod m\) khi và chỉ
khi \(\phi(m)\mid n\). Hơn nữa, \(g^{n_0}\equiv g^{n_1}\) khi và chỉ khi
\(\phi(m)\mid(n_0-n_1)\).

\paragraph{Định lý 8.6.}
Nếu \(g\) là căn nguyên thủy của \(m\), thì các lũy thừa
\[
 1,g,g^2,\ldots,g^{\phi(m)-1}
\]
đại diện duy nhất cho các số nguyên tố cùng nhau với \(m\) modulo \(m\). Đặc
biệt, nếu \(m>2\), thì
\[
 g^{\phi(m)/2}\equiv-1\pmod m.
\]

\paragraph{Chứng minh.}
Các lũy thừa đều nguyên tố cùng nhau với \(m\), và có đúng \(\phi(m)\) số như
vậy. Nếu \(g^i\equiv g^j\), thì \(g^{i-j}\equiv1\), nên \(\phi(m)\mid(i-j)\).
Do đó chúng phân biệt và phủ hết các lớp khả nghịch. Với \(-1\), tồn tại duy
nhất \(i\) sao cho \(g^i\equiv-1\); bình phương cho \(g^{2i}\equiv1\), nên
\(2i=\phi(m)\) trong khoảng \(0\le i\le\phi(m)-1\).

\paragraph{Mệnh đề 8.7.}
Với số nguyên dương \(m\), đồng dư \(x^2\equiv1\pmod m\) chỉ có các nghiệm
\(x\equiv\pm1\pmod m\) khi và chỉ khi \(m\) có căn nguyên thủy.

\paragraph{Ví dụ 8.2.}
Với số nguyên dương \(m\), gọi \(S\) là tập các số nguyên dương nhỏ hơn \(m\)
và nguyên tố cùng nhau với \(m\), và \(P\) là tích các phần tử của \(S\). Chứng
minh \(P\equiv\pm1\pmod m\), với \(P\equiv-1\pmod m\) khi và chỉ khi \(m\) có
căn nguyên thủy.

\paragraph{Lời giải.}
Với \(m\ge3\), chia \(S\) thành \(A\), các nghiệm của \(x^2\equiv1\pmod m\),
và \(B\), các phần tử còn lại. Các phần tử của \(B\) ghép cặp với nghịch đảo
nhân riêng biệt, nên tích của chúng là \(1\). Các phần tử của \(A\) ghép cặp
với nghịch đảo cộng \(x\) và \(m-x\); mỗi cặp có tích \(-1\). Nếu \(m\) có căn
nguyên thủy, theo Mệnh đề 8.7, \(A=\{1,-1\}\), nên \(P\equiv-1\). Nếu không,
theo Ví dụ 4.9, số phần tử của \(A\) là một lũy thừa của \(2\) ít nhất \(4\),
nên số cặp là chẵn và \(P\equiv1\). Với \(m\) nguyên tố, đây chính là định lý
Wilson.

\paragraph{Định lý 8.8.}
\begin{enumerate}
\item Nếu \(g\) là căn nguyên thủy của số nguyên tố \(p\), thì \(g\) hoặc
\(g+p\) là căn nguyên thủy của \(p^2\), tùy theo
\[
 g^{p-1}\not\equiv1\pmod{p^2}
 \quad\text{hoặc}\quad
 g^{p-1}\equiv1\pmod{p^2}.
\]
\item Nếu \(g\) là căn nguyên thủy của \(p^k\), với \(k\ge2\) và \(p\) nguyên
tố, thì \(g\) là căn nguyên thủy của \(p^{k+1}\).
\end{enumerate}

Theo Định lý 8.6, nếu biết căn nguyên thủy \(g\) của \(m\), thì với mỗi \(a\)
nguyên tố cùng nhau với \(m\), tồn tại duy nhất \(i\) modulo \(\phi(m)\) sao
cho \(g^i\equiv a\pmod m\). Số \(i\) gọi là chỉ số của \(a\) theo cơ số \(g\),
ký hiệu \(\ind_g(a)\). Chỉ số giống logarit qua các tính chất:
\[
\begin{gathered}
\ind_g(1)\equiv0,\qquad \ind_g(g)\equiv1\pmod{\phi(m)},\\
a\equiv b\pmod m\Rightarrow \ind_g(a)\equiv\ind_g(b)\pmod{\phi(m)},\\
\ind_g(ab)\equiv\ind_g(a)+\ind_g(b)\pmod{\phi(m)},\\
\ind_g(a^k)\equiv k\ind_g(a)\pmod{\phi(m)}.
\end{gathered}
\]

\paragraph{Định lý 8.9.}
Nếu \(p\) là số nguyên tố và \(p\nmid a\), thì đồng dư \(x^n\equiv a\pmod p\)
có \(\gcd(n,p-1)\) nghiệm hoặc không có nghiệm tùy theo
\[
 a^{(p-1)/\gcd(n,p-1)}\equiv1\pmod p
\]
hoặc không.

\paragraph{Chứng minh.}
Lấy căn nguyên thủy \(g\) của \(p\), gọi \(i=\ind_g(a)\) và \(u=\ind_g(x)\).
Đồng dư trở thành
\[
 nu\equiv i\pmod{p-1}.
\]
Đặt \(k=\gcd(n,p-1)\). Đồng dư tuyến tính này có \(k\) nghiệm khi và chỉ khi
\(k\mid i\). Vì \(g\) là căn nguyên thủy, điều này tương đương
\[
 g^{i(p-1)/k}\equiv a^{(p-1)/k}\equiv1\pmod p.
\]

\paragraph{Ví dụ 8.3.}
Cho \(n\ge2\) và \(p=2^n+1\). Chứng minh rằng nếu
\[
 3^{(p-1)/2}+1\equiv0\pmod p,
\]
thì \(p\) là số nguyên tố.

\paragraph{Lời giải.}
Từ \(3^{2^{n-1}}\equiv-1\pmod p\), ta được \(3^{2^n}\equiv1\pmod p\), nên cấp
của \(3\) modulo \(p\) là \(2^n=p-1\). Cấp luôn chia \(\phi(p)\le p-1\), do đó
\(\phi(p)=p-1\), suy ra \(p\) nguyên tố.

\paragraph{Ví dụ 8.4.}
Chứng minh rằng nếu \(n=3^{k-1}\), thì \(2^n\equiv-1\pmod{3^k}\).

\paragraph{Lời giải.}
Theo Ví dụ 8.1, \(2\) là căn nguyên thủy của \(3^k\), nên cấp của \(2\) là
\(\phi(3^k)=2\cdot3^{k-1}=2n\). Do đó \(2^{2n}\equiv1\). Nhưng
\[
 2^n-1\equiv(-1)^{3^{k-1}}-1\equiv1\not\equiv0\pmod3,
\]
nên \(2^n+1\equiv0\pmod{3^k}\).

\paragraph{Ví dụ 8.5.}
Tìm mọi số nguyên dương \(n>1\) sao cho
\[
 \frac{2^n+1}{n^2}
\]
là số nguyên. (IMO 1990)

\paragraph{Lời giải.}
Rõ ràng \(n\) lẻ. Giả sử \(3^k\|n\). Khi đó
\[
 3^{2k}\mid n^2\mid(2^n+1),
\]
nên \(2^n\equiv-1\pmod{3^{2k}}\). Theo Hệ quả 8.2,
\(3^{2k-1}\mid n\), suy ra \(2k-1\le k\), tức \(k\le1\). Ta thấy \(n=3\) là
nghiệm.

Giả sử \(n\) có ước nguyên tố \(p>3\), lấy \(p\) nhỏ nhất. Từ
\(p\mid2^n+1\), \(2^n\equiv-1\pmod p\). Gọi \(d\) là cấp của \(2\) modulo
\(p\). Vì \(2^{2n}\equiv1\), \(d\mid2n\). Nếu \(d\) lẻ thì \(d\mid n\), mâu
thuẫn với \(2^n\equiv-1\); vậy \(d=2d_1\). Khi đó \(d_1\mid n\), và
\(2d_1\mid p-1\), nên \(d_1<(p)\). Do tính nhỏ nhất của \(p\), \(d_1=1\) hoặc
\(3\). Nếu \(d_1=1\), thì \(d=2\), vô lý. Nếu \(d_1=3\), thì \(d=6\) và
\(p\mid63\), nên \(p=7\), nhưng cấp của \(2\) modulo \(7\) là \(3\), mâu
thuẫn. Vậy không có \(p>3\); nghiệm duy nhất là \(n=3\).

\paragraph{Sự kiện hữu ích.}
Mọi ước nguyên tố của số Fermat \(2^{2^n}+1\), \(n>1\), đều có dạng
\[
 2^{n+2}k+1.
\]

\paragraph{Bài tập.}
\begin{enumerate}
\item Cho \(p\) là số nguyên tố lẻ. Chứng minh
\[
 1^i+2^i+\cdots+(p-1)^i\equiv0\pmod p
\]
với mọi \(i\), \(0\le i\le p-2\).
\item Chứng minh rằng nếu \(p\) là số nguyên tố lẻ, thì \(x^4\equiv-1\pmod p\)
có nghiệm khi và chỉ khi \(p\equiv1\pmod8\).
\item Chứng minh rằng nếu \(a,n\) là các số nguyên dương và \(a\) lẻ, thì
\[
 a^{2^n}\equiv1\pmod{2^{n+2}}.
\]
\item Số \(142857\) có tính chất đáng chú ý: nhân nó với \(1,2,3,4,5,6\) sẽ
chỉ hoán vị vòng các chữ số. Còn những số nào có tính chất này? Gợi ý: tính
\(142857\times7\).
\end{enumerate}

\section{Chuỗi Dirichlet}

Dù tên gọi có vẻ đáng ngại, chuỗi Dirichlet dễ thao tác và có thể cho các
chứng minh nhanh cho một số đẳng thức số học, chẳng hạn Ví dụ 3.2. Cho
\(\alpha\) là hàm từ các số nguyên dương vào các số nguyên. Ta gọi
\[
 f(s)=\sum_{n=1}^{\infty}\frac{\alpha(n)}{n^s}
 =\alpha(1)+\frac{\alpha(2)}{2^s}+\frac{\alpha(3)}{3^s}+\cdots
\]
là hàm sinh chuỗi Dirichlet của \(\alpha\), ký hiệu \(f(s)\leftrightarrow
\alpha(n)\). Như các hàm sinh nói chung, chúng được dùng để lấy thông tin về
các hàm số học tương ứng qua việc biến đổi hàm sinh.

Gọi \(1\) là hàm luôn bằng \(1\) trên các số nguyên dương. Định nghĩa
\[
\delta_1(n)=
\begin{cases}
1,&n=1,\\
0,&n>1.
\end{cases}
\]
Cả \(1\) và \(\delta_1\) đều nhân tính.

Với hai hàm \(\alpha,\beta\), tích chập của chúng, ký hiệu \(\alpha*\beta\),
được định nghĩa bởi
\[
 (\alpha*\beta)(n)=\sum_{d\mid n}\alpha(d)\beta(n/d).
\]
Tích chập đối xứng: \(\alpha*\beta=\beta*\alpha\).

\paragraph{Định lý 9.1.}
Nếu \(f(s)\leftrightarrow\alpha(n)\) và \(g(s)\leftrightarrow\beta(n)\), thì
\[
 (fg)(s)\leftrightarrow(\alpha*\beta)(n).
\]

Hàm sinh của \(1(n)\) là hàm zeta Riemann quen thuộc:
\[
 \zeta(s)=\sum_{n=1}^{\infty}\frac1{n^s}=1+\frac1{2^s}+\frac1{3^s}+\cdots,
\]
nên \(\zeta(s)\leftrightarrow1(n)\). Ta làm việc với các hàm này ở mức hình
thức; không cần biết tính chất giải tích của \(\zeta(s)\) hay của chuỗi
Dirichlet nào khác. Theo Định lý 9.1,
\[
 \zeta^2(s)\leftrightarrow \sum_{d\mid n}1=\tau(n).
\]
Cuối cùng, rõ ràng \(1\leftrightarrow\delta_1(n)\).

\paragraph{Định lý 9.2.}
Nếu \(\alpha\) là hàm nhân tính, thì
\[
 \sum_{n=1}^{\infty}\frac{\alpha(n)}{n^s}
 =
 \prod_p\sum_{k=0}^{\infty}\frac{\alpha(p^k)}{p^{ks}}
 =
 \prod_p\left(1+\alpha(p)p^{-s}+\alpha(p^2)p^{-2s}+\cdots\right),
\]
trong đó tích chạy trên mọi số nguyên tố \(p\).

Lấy \(\alpha=1\), ta được
\[
 \zeta(s)=\prod_p(1+p^{-s}+p^{-2s}+\cdots)
 =\prod_p\frac1{1-p^{-s}}.
\]

Số nguyên dương \(n>1\) gọi là không chứa bình phương nếu không có số nguyên
tố \(p\) nào thỏa \(p^2\mid n\). Định nghĩa hàm Möbius \(\mu\) bởi
\[
\mu(n)=
\begin{cases}
1,&n=1,\\
0,&n\text{ không không chứa bình phương},\\
(-1)^k,&n\text{ không chứa bình phương và có }k\text{ ước nguyên tố}.
\end{cases}
\]
Hàm \(\mu\) nhân tính. Theo Định lý 9.2,
\[
 \prod_p(1-p^{-s})=\frac1{\zeta(s)}.
\]
Vậy \(1/\zeta(s)\leftrightarrow\mu(n)\). Tính chất này làm hàm \(\mu\), dù
thoạt nhìn bí ẩn, trở nên rất quan trọng.

\paragraph{Định lý 9.3 (công thức đảo Möbius).}
Cho \(\alpha,\beta\) là các hàm sao cho
\[
 \beta(n)=\sum_{d\mid n}\alpha(d).
\]
Khi đó
\[
 \alpha(n)=\sum_{d\mid n}\mu(n/d)\beta(d).
\]

\paragraph{Chứng minh.}
Gọi \(f(s)\leftrightarrow\alpha(n)\), \(g(s)\leftrightarrow\beta(n)\). Điều
kiện tương đương \(\beta=\alpha*1\), tức \(g(s)=f(s)\zeta(s)\). Kết luận tương
đương \(\alpha=\beta*\mu\), tức \(f(s)=g(s)/\zeta(s)\).

\paragraph{Định lý 9.4.}
Nếu \(f(s)\leftrightarrow\alpha(n)\), thì với mọi số nguyên \(k\),
\[
 f(s-k)\leftrightarrow n^k\alpha(n).
\]

Để đọc thêm về chuỗi Dirichlet và hàm sinh nói chung, xem H. Wilf,
\textit{Generatingfunctionology}.

\paragraph{Bài tập.}
\begin{enumerate}
\item Cho \(\alpha,\beta,\gamma\) là các hàm từ số nguyên dương vào số nguyên.
\begin{enumerate}[label=(\alph*)]
\item Chứng minh \(\alpha*\delta_1=\alpha\).
\item Chứng minh \((\alpha*\beta)*\gamma=\alpha*(\beta*\gamma)\).
\item Chứng minh nếu \(\alpha,\beta\) nhân tính, thì \(\alpha*\beta\) cũng
nhân tính.
\end{enumerate}
\item Chứng minh các quan hệ
\[
 \frac{\zeta(s-1)}{\zeta(s)}\leftrightarrow\phi(n),\qquad
 \zeta(s)\zeta(s-1)\leftrightarrow\sigma(n),\qquad
 \frac{\zeta(s)}{\zeta(2s)}\leftrightarrow|\mu(n)|.
\]
\item Cho phân tích nguyên tố của \(n>1\) là \(p_1^{e_1}\cdots p_k^{e_k}\).
Định nghĩa \(\lambda(n)=(-1)^{e_1+\cdots+e_k}\) và \(\theta(n)=2^k\), với
\(\lambda(1)=\theta(1)=1\). Chứng minh \(\lambda,\theta\) nhân tính và
\[
 \frac{\zeta(2s)}{\zeta(s)}\leftrightarrow\lambda(n),\qquad
 \frac{\zeta^2(s)}{\zeta(2s)}\leftrightarrow\theta(n).
\]
\item Với mọi số nguyên dương \(n\), đặt
\[
 f(n)=\sum_{m=1}^n\frac n{\gcd(m,n)}.
\]
\begin{enumerate}[label=(\alph*)]
\item Chứng minh \(f(n)=\sum_{d\mid n}d\phi(d)\).
\item Nếu \(n=p_1^{e_1}\cdots p_k^{e_k}>1\), chứng minh
\[
 f(n)=\frac{p_1^{2e_1+1}+1}{p_1+1}
 \frac{p_2^{2e_2+1}+1}{p_2+1}\cdots
 \frac{p_k^{2e_k+1}+1}{p_k+1}.
\]
\end{enumerate}
\item Kiểm tra Ví dụ 3.2 bằng một phép tính.
\item Gọi \(\operatorname{id}\) là hàm đồng nhất, tức \(\operatorname{id}(n)=n\).
Kiểm tra trong một phép tính từng đẳng thức:
\[
 \phi*\tau=\sigma,\quad
 \mu*1=\delta_1,\quad
 \mu*\operatorname{id}=\phi,\quad
 \phi*\sigma=\operatorname{id}\cdot\tau,\quad
 \sigma*\operatorname{id}=1*(\operatorname{id}\cdot\tau).
\]
\item Cho \(a_1,a_2,\ldots\) là dãy số nguyên dương thỏa
\[
 \sum_{d\mid n}a_d=2^n
\]
với mọi \(n\). Do đó \(a_1=2\), \(a_2=2^2-2=2\),
\(a_3=2^3-2=6\), \(a_4=2^4-2-2=12\), v.v. Chứng minh \(n\mid a_n\) với mọi
\(n\). Gợi ý: đừng dùng hàm sinh chuỗi Dirichlet của \((a_n)\); dùng công thức
đảo Möbius. Gợi ý lớn hơn: xét hàm \(f:[0,1]\to[0,1]\) cho bởi \(f(x)=\{2x\}\),
trong đó \(\{x\}=x-\lfloor x\rfloor\). Tìm liên hệ giữa công thức trong bài và
hàm \(f^{(n)}=f\circ f\circ\cdots\circ f\).
\item Với mọi số nguyên không âm \(k\), định nghĩa
\[
 \sigma_k(n)=\sum_{d\mid n}d^k.
\]
Như vậy \(\sigma_0=\tau\) và \(\sigma_1=\sigma\). Chứng minh
\[
 \zeta(s)\zeta(s-k)\leftrightarrow\sigma_k(n).
\]
\end{enumerate}

\section{Một số chủ đề khác}

\subsection{Phương trình Pell}

Phương trình Pell, hay đúng hơn là phương trình Fermat, là các phương trình
Diophantine dạng
\[
 x^2-dy^2=N,
\]
trong đó \(d\) là số nguyên dương không chính phương. Khi \(N=1\), luôn có vô
hạn nghiệm; ta mô tả chúng.

\paragraph{Định lý 10.1.1.}
Nếu \((a,b)\) là nghiệm nguyên dương nhỏ nhất của \(x^2-dy^2=1\), thì mọi
nghiệm nguyên dương có dạng
\[
 (x_n,y_n)=\left(
 \frac{(a+b\sqrt d)^n+(a-b\sqrt d)^n}{2},
 \frac{(a+b\sqrt d)^n-(a-b\sqrt d)^n}{2\sqrt d}
 \right).
\]

Ta không chứng minh đầy đủ ở đây, nhưng kiểm tra rằng mọi cặp trong công thức
đều là nghiệm. Chúng thỏa
\[
 x_n+y_n\sqrt d=(a+b\sqrt d)^n,\qquad
 x_n-y_n\sqrt d=(a-b\sqrt d)^n.
\]
Do đó
\[
 x_n^2-dy_n^2=(a+b\sqrt d)^n(a-b\sqrt d)^n=(a^2-db^2)^n=1.
\]
Hai dãy \((x_n)\), \((y_n)\) thỏa truy hồi
\[
 x_n=2ax_{n-1}-x_{n-2},\qquad y_n=2ay_{n-1}-y_{n-2}.
\]
Với \(x^2-dy^2=-1\), tình hình tương tự: nếu \((a,b)\) là nghiệm dương nhỏ
nhất, thì các \((x_n,y_n)\) ở trên với \(n\) lẻ là nghiệm của
\(x^2-dy^2=-1\), còn với \(n\) chẵn là nghiệm của \(x^2-dy^2=1\).

\paragraph{Ví dụ 10.1.1.}
Tìm mọi nghiệm nguyên dương \((x,y)\) của \(x^2-2y^2=1\).

\paragraph{Lời giải.}
Nghiệm nguyên dương nhỏ nhất là \((3,2)\), nên mọi nghiệm là
\[
 (x_n,y_n)=\left(
 \frac{(3+2\sqrt2)^n+(3-2\sqrt2)^n}{2},
 \frac{(3+2\sqrt2)^n-(3-2\sqrt2)^n}{2\sqrt2}
 \right).
\]
Các nghiệm đầu là \((3,2)\), \((17,12)\), \((99,70)\).

\paragraph{Ví dụ 10.1.2.}
Chứng minh phương trình \(x^2-dy^2=-1\) không có nghiệm nguyên nếu
\(d\equiv3\pmod4\).

\paragraph{Lời giải.}
Số \(d\) có một ước nguyên tố dạng \(4k+3\), gọi là \(q\). Khi đó
\(x^2\equiv-1\pmod q\), mâu thuẫn với Định lý 4.9.

\paragraph{Bài tập.}
\begin{enumerate}
\item Trong dãy
\[
 \frac12,\frac53,\frac{11}8,\frac{27}{19},\ldots,
\]
mẫu của hạng thứ \(n\) (\(n>1\)) là tổng tử và mẫu của hạng thứ \(n-1\). Tử
của hạng thứ \(n\) là tổng các mẫu của hạng thứ \(n\) và \(n-1\). Tìm giới hạn
của dãy. (ARML 1979)
\item Cho \(x_0=0,x_1=1,x_{n+1}=4x_n-x_{n-1}\), và
\(y_0=1,y_1=2,y_{n+1}=4y_n-y_{n-1}\). Chứng minh với mọi \(n\ge0\),
\[
 y_n^2=3x_n^2+1.
\]
(Olympiad Toán Canada 1988)
\item Các đa thức \(P,Q\) có \(\deg P=n\), \(\deg Q=m\), cùng hệ số đầu, và
\[
 P^2(x)=(x^2-1)Q^2(x)+1.
\]
Chứng minh \(P'(x)=nQ(x)\). (Chung kết Olympiad Toán Thụy Điển 1978)
\end{enumerate}

\subsection{Dãy Farey}

Dãy Farey bậc \(n\) là dãy tất cả phân số tối giản trong \([0,1]\), với cả tử
và mẫu không vượt quá \(n\), sắp theo thứ tự tăng. Năm dãy Farey đầu là
\[
\begin{array}{lllllllllll}
0/1,&&&&&&&&&&1/1,\\
0/1,&&&&1/2,&&&&&&1/1,\\
0/1,&&1/3,&&1/2,&&2/3,&&&&1/1,\\
0/1,&1/4,&1/3,&&1/2,&&2/3,&3/4,&&&1/1,\\
0/1,&1/5,&1/4,&1/3,&2/5,&1/2,&3/5,&2/3,&3/4,&4/5,&1/1.
\end{array}
\]

Một số tính chất:
\begin{enumerate}
\item Nếu \(a/b\) và \(c/d\) là hai phân số liên tiếp trong cùng một dãy, theo
thứ tự đó, thì \(ad-bc=1\).
\item Nếu \(a/b,c/d,e/f\) là ba phân số liên tiếp trong cùng một dãy, theo thứ
tự đó, thì
\[
 \frac{a+e}{b+f}=\frac cd.
\]
\item Nếu \(a/b\) và \(c/d\) liên tiếp trong cùng một dãy, thì trong mọi phân
số nằm giữa chúng, \((a+c)/(b+d)\) sau khi đưa về dạng tối giản là phân số duy nhất có mẫu
nhỏ nhất.
\end{enumerate}
Để xem chứng minh và các tính chất thú vị khác, có thể đọc Ross Honsberger,
\textit{Farey Sequences}, trong \textit{Ingenuity in Mathematics}.

\paragraph{Bài tập.}
\begin{enumerate}
\item Gọi \(a_1,a_2,\ldots,a_m\) là các mẫu số của các phân số trong dãy Farey
bậc \(n\), theo thứ tự. Chứng minh
\[
 \frac1{a_1a_2}+\frac1{a_2a_3}+\cdots+\frac1{a_{m-1}a_m}=1.
\]
\end{enumerate}

\subsection{Phân số liên tục}

Cho \(a_0,a_1,\ldots,a_n\) là các số thực, tất cả dương trừ có thể \(a_0\).
Ký hiệu \(\langle a_0,a_1,\ldots,a_n\rangle\) là phân số liên tục
\[
 a_0+\frac1{a_1+\frac1{\ddots+\frac1{a_{n-1}+\frac1{a_n}}}}.
\]
Nếu mọi \(a_i\) là số nguyên, phân số liên tục gọi là đơn. Định nghĩa các dãy
\((p_k)\), \((q_k)\) bởi
\[
 p_{-1}=0,\quad p_0=a_0,\quad p_k=a_kp_{k-1}+p_{k-2},
\]
\[
 q_{-1}=0,\quad q_0=1,\quad q_k=a_kq_{k-1}+q_{k-2}\qquad(k\ge1).
\]

\paragraph{Định lý 10.3.1.}
Với mọi \(x>0\) và \(k\ge1\),
\[
 \langle a_0,a_1,\ldots,a_{k-1},x\rangle
 =
 \frac{xp_{k-1}+p_{k-2}}{xq_{k-1}+q_{k-2}}.
\]
Đặc biệt,
\[
 \langle a_0,a_1,\ldots,a_k\rangle=\frac{p_k}{q_k}.
\]

\paragraph{Định lý 10.3.2.}
Với mọi \(k\ge0\),
\[
 p_kq_{k-1}-p_{k-1}q_k=(-1)^{k-1},
\]
\[
 p_kq_{k-2}-p_{k-2}q_k=(-1)^ka_k.
\]

Gọi \(c_k\) là hội tụ thứ \(k\):
\[
 c_k=\langle a_0,a_1,\ldots,a_k\rangle=\frac{p_k}{q_k}.
\]

\paragraph{Định lý 10.3.3.}
\[
 c_0<c_2<c_4<\cdots<c_5<c_3<c_1.
\]
Một liên hệ đẹp giữa phân số liên tục, phương trình Diophantine tuyến tính và
phương trình Pell có trong Andy Liu, \textit{Continued Fractions and
Diophantine Equations}, Mathematical Mayhem, tập 3, số 2.

\paragraph{Bài tập.}
\begin{enumerate}
\item Cho \(a=\langle1,2,\ldots,99\rangle\) và
\(b=\langle1,2,\ldots,99,100\rangle\). Chứng minh
\[
 |a-b|<\frac1{99!\,100!}.
\]
(Tournament of Towns 1990)
\item Tính
\[
 \sqrt[8]{2207-\frac1{2207-\frac1{2207-\cdots}}}.
\]
Viết đáp án dưới dạng \(\frac{a+b\sqrt c}{d}\), trong đó \(a,b,c,d\) là các
số nguyên. (Putnam 1995)
\end{enumerate}

\subsection{Bài toán tem thư}

Cho \(a,b>1\) là hai số nguyên dương nguyên tố cùng nhau. Xét tập các số có
dạng \(ax+by\), trong đó \(x,y\) là các số nguyên không âm. Khi đó:
\begin{enumerate}
\item Số nguyên lớn nhất không viết được dưới dạng trên là
\[
 (a-1)(b-1)-1=ab-a-b.
\]
\item Có \(\frac12(a-1)(b-1)\) số nguyên dương không viết được dưới dạng trên.
\item Với mọi số nguyên \(t\), \(0\le t\le ab-a-b\), \(t\) viết được dưới dạng
trên khi và chỉ khi \(ab-a-b-t\) không viết được.
\end{enumerate}
Nếu chưa từng thử bài toán hấp dẫn này, nên tự thử trước khi đọc lời giải đầy
đủ.

Trước lời giải, xét ví dụ \(a=12\), \(b=5\). Các số không âm đầu tiên, xếp
thành các hàng độ dài \(12\), trong đó các số không biểu diễn được được đặt
trong ngoặc, là
\[
\begin{array}{rrrrrrrrrrrr}
0&(1)&(2)&(3)&(4)&5&(6)&(7)&(8)&(9)&10&(11)\\
12&(13)&(14)&15&(16)&17&18&(19)&20&21&22&(23)\\
24&25&(26)&27&28&29&30&31&32&33&34&35\\
36&37&38&39&40&41&42&43&44&45&46&47\\
48&49&50&51&52&53&54&55&56&57&58&59
\end{array}
\]
Quan sát chính là trong mỗi cột, các số không biểu diễn được kết thúc khi ta
gặp một bội của \(5\). Đi xuống một cột là cộng \(12\); vì thế một khi gặp một
số biểu diễn được, mọi số bên dưới cũng biểu diễn được. Đây là chìa khóa của
lời giải.

\paragraph{Chứng minh.}
Gọi một số là biểu diễn được nếu nó có dạng \(ax+by\) với \(x,y\ge0\). Với
mỗi \(i\), \(0\le i\le a-1\), gọi \(m_i\) là số nguyên không âm nhỏ nhất sao
cho
\[
 b\mid(i+am_i).
\]
Rõ ràng với \(k\ge m_i\), số \(i+ak\) biểu diễn được. Ta khẳng định với
\(0\le k\le m_i-1\), \(i+ak\) không biểu diễn được; chỉ cần chứng minh
\(i+a(m_i-1)\) không biểu diễn được nếu \(m_i\ge1\).

Viết \(i+am_i=bn_i\). Vì
\[
 i+a(m_i-b)=b(n_i-a),
\]
nên \(m_i<b\), nếu không sẽ có \(m_i\) nhỏ hơn. Khi đó
\[
 i+a(m_i-b)\le a-1-a=-1,
\]
suy ra \(n_i<a\). Giả sử
\[
 ax+by=i+a(m_i-1)=bn_i-a.
\]
Khi đó \(a(x+1)=b(n_i-y)\), nên \(n_i-y>0\). Vì \(a,b\) nguyên tố cùng nhau,
\(a\mid n_i-y\), nhưng \(0<n_i-y\le n_i\le a-1\), mâu thuẫn.

Vậy số không biểu diễn được lớn nhất có dạng \(bn_i-a\) với \(n_i\le a-1\),
và lập luận trên cũng cho thấy mọi số dương dạng này đều không biểu diễn
được. Số lớn nhất là \(b(a-1)-a=ab-a-b\), chứng minh (1).

Trong cột \(i\) có \(m_i\) số không biểu diễn được. Số biểu diễn được đầu tiên
trong cột \(i\) là \(n_ib\). Vì \(0,b,2b,\ldots,(a-1)b\) nằm ở các cột khác
nhau, nên khi \(i\) chạy từ \(0\) đến \(a-1\), các \(n_i\) nhận đúng các giá
trị \(0,1,\ldots,a-1\). Do đó
\[
\begin{aligned}
\sum_i(i+am_i)
&=\sum_i i+a\sum_i m_i\\
&=\sum_i bn_i=\frac{a(a-1)b}{2}.
\end{aligned}
\]
Mà \(\sum_i i=a(a-1)/2\), nên
\[
 \sum_i m_i=\frac{(a-1)(b-1)}2,
\]
chứng minh (2).

Cuối cùng, nếu cả \(t\) và \(ab-a-b-t\) đều biểu diễn được, tức
\[
 ax_1+by_1=t,\qquad ax_2+by_2=ab-a-b-t,
\]
thì
\[
 a(x_1+x_2)+b(y_1+y_2)=ab-a-b,
\]
mâu thuẫn với (1). Xét các cặp \((t,ab-a-b-t)\) với
\[
 0\le t\le\frac{(a-1)(b-1)}2-1,
\]
mỗi cặp có nhiều nhất một số biểu diễn được. Nhưng ta đã biết có đúng
\((a-1)(b-1)/2\) số không biểu diễn được, bằng số cặp, nên mỗi cặp có đúng một
số biểu diễn được. Điều này chứng minh (3).

Có một chứng minh ngắn hơn nhiều bằng Hệ quả 2.4; bạn có tìm được không?

Kiểu bài toán này thể hiện rõ tinh thần giải toán trong lý thuyết số và trong
toán học nói chung. Nếu chỉ trình bày chứng minh, nó có thể trông nhân tạo và
thiếu động cơ. Nhưng bằng cách xét một ví dụ cụ thể, tìm mẫu hình, rồi dùng
mẫu hình đó làm điểm tựa, ta mở rộng thành chứng minh đầy đủ. Phần đại số
trong chứng minh thực chất chỉ là dịch các mẫu hình quan sát được sang ký hiệu
hình thức. Toán học có thể được mô tả như việc nghiên cứu mẫu hình. Cũng đáng
chú ý rằng ta chỉ dùng các kết quả rất sơ cấp, cho thấy các khái niệm cơ bản
có sức mạnh lớn đến mức nào. Lời giải có thể hơi nhiều chi tiết, nhưng đừng
ngại làm việc với chi tiết; càng đi sâu, ta càng hiểu tầm quan trọng của các
khái niệm này và các quan hệ tinh tế giữa chúng. Khi theo đuổi một ý tưởng
đến cùng, đôi khi ta cảm thấy chứng minh gần như tự vận hành. Bài học là: một
ý tưởng đơn giản có thể đi rất xa.

Để đọc thêm về bài toán tem thư, xem Ross Honsberger, \textit{A Putnam Paper
Problem}, trong \textit{Mathematical Gems II}.

\paragraph{Bài tập.}
\begin{enumerate}
\item Cho \(a,b,c\) là các số nguyên dương đôi một không có ước chung lớn hơn
\(1\). Chứng minh rằng
\[
 2abc-ab-bc-ca
\]
là số nguyên lớn nhất không biểu diễn được dưới dạng
\[
 xab+yca+zab,
\]
trong đó \(x,y,z\) là các số nguyên không âm. (IMO 1983)
\end{enumerate}

\section*{Tài liệu tham khảo}

A. Adler và J. Coury, \textit{The Theory of Numbers}, Jones and Bartlett.

I. Niven và H. Zuckerman, \textit{An Introduction to the Theory of Numbers},
John Wiley \& Sons.

Phiên bản thứ nhất: tháng 10 năm 1995. Phiên bản thứ hai: tháng 1 năm 1996.
Phiên bản thứ ba: tháng 4 năm 1999. Phiên bản thứ tư: tháng 5 năm 2000.

Tác giả cảm ơn Ather Gattami vì một cải tiến trong chứng minh bài toán tem
thư. Tài liệu này được dàn trang bằng \LaTeX{} và có thể được phân phối tự do
miễn là nội dung không bị thay đổi và thông báo bản quyền không bị xóa. Mọi
góp ý và chỉnh sửa đều được hoan nghênh. Tài liệu không được bán để thu lợi
hoặc đưa vào các tài liệu thương mại nếu không có sự cho phép rõ ràng của chủ
sở hữu bản quyền.

\end{document}
